% =====================================================================
%  CAPITULO 10 - CAPACITORES E INDUTORES
% =====================================================================
\section{Capacitores e Indutores}

\begin{conceito}[Antes de começar]
\textbf{Capacitor plano:} $C = \dfrac{\varepsilon_0\,A}{d}$, com
$\varepsilon_0 = \SI{8.85e-12}{\farad\per\meter}$.\par
\textbf{Carga e energia:} $Q = C\,U$ \quad e \quad
$E = \tfrac12 C\,U^2 = \tfrac12 Q\,U$.\par
\textbf{Associação (oposta à dos resistores):}
\begin{itemize}[leftmargin=1.6em,itemsep=1pt,topsep=2pt]
\item \emph{paralelo:} $\Ceq = C_1 + C_2 + \dots$ (as capacitâncias somam);
\item \emph{série:} $\dfrac{1}{\Ceq} = \dfrac{1}{C_1}+\dfrac{1}{C_2}+\dots$
\end{itemize}
\textbf{Circuito RC --- constante de tempo:} $\tau = R\,C$. Na \emph{carga},
$v_C(t) = U_0\big(1-e^{-t/\tau}\big)$; na \emph{descarga},
$v_C(t) = U_0\,e^{-t/\tau}$. Após $5\tau$, considera-se o regime permanente
atingido.\par
\textbf{Indutor:} $E = \tfrac12 L\,I^2$; a associação segue a \emph{mesma} regra
dos resistores (série soma, paralelo inverso). Constante de tempo RL:
$\tau = L/R$.
\end{conceito}

\legendaniveis

% =====================================================================
\subsection{Capacitores}
% =====================================================================

\blocofixacao

\begin{questao}[1]{}
A capacitância de um capacitor de placas paralelas:
\begin{alternativas}
\item aumenta com a distância entre as placas.
\item aumenta com a área das placas.
\item independe do meio entre as placas.
\item diminui ao inserir um dielétrico.
\item depende da tensão aplicada.
\end{alternativas}
\resposta{(b)}
\resolucao{De $C = \dfrac{\varepsilon_0 A}{d}$: a capacitância cresce com a área
$A$ e \emph{diminui} com a distância $d$ --- o que elimina (a). Ela não depende da
tensão (elimina e), e inserir um dielétrico \emph{aumenta} $C$ (elimina d).}
\end{questao}

\begin{questao}[1]{}
Calcule a capacitância de um capacitor plano de área $A = \SI{100}{\centi\meter\squared}$
e distância entre placas $d = \SI{1}{\milli\meter}$, no vácuo.
\dados{$\varepsilon_0 = \SI{8.85e-12}{\farad\per\meter}$.}
\resposta{$C \approx \SI{88.5}{\pico\farad}$}
\resolucao{Convertendo: $A = \SI{100}{\centi\meter\squared} = \pot{100}{-4}\,\si{\meter\squared}
= \pot{1}{-2}\,\si{\meter\squared}$ e $d = \pot{1}{-3}\,\si{\meter}$:
\[
C=\frac{\varepsilon_0 A}{d}
=\frac{\pot{8.85}{-12}\cdot\pot{1}{-2}}{\pot{1}{-3}}
=\pot{8.85}{-11}\,\si{\farad}\approx\SI{88.5}{\pico\farad}.
\]}
\end{questao}

\begin{questao}[1]{}
Um capacitor de $\SI{10}{\micro\farad}$ é submetido a uma tensão de
$\SI{12}{\volt}$. Determine a carga armazenada e a energia acumulada.
\resposta{$Q = \SI{120}{\micro\coulomb}$; $E = \SI{720}{\micro\joule}$}
\resolucao{$Q = C\,U = \pot{10}{-6}\cdot12 = \SI{120}{\micro\coulomb}$.\\
$E = \tfrac12 C\,U^2 = \tfrac12\cdot\pot{10}{-6}\cdot12^2
= \SI{720}{\micro\joule}$.}
\end{questao}

\begin{questao}[1]{}
Dois capacitores de $\SI{6}{\micro\farad}$ e $\SI{3}{\micro\farad}$ são ligados
em \textbf{série}. A capacitância equivalente é:
\begin{alternativas}[3]
\item \SI{2}{\micro\farad}
\item \SI{4.5}{\micro\farad}
\item \SI{9}{\micro\farad}
\item \SI{18}{\micro\farad}
\item \SI{0.5}{\micro\farad}
\end{alternativas}
\resposta{(a)}
\resolucao{Em série, $\Ceq = \dfrac{C_1 C_2}{C_1 + C_2} = \dfrac{6\cdot3}{9}
= \SI{2}{\micro\farad}$. Note que, no capacitor, a série se comporta como o
paralelo do resistor: o equivalente é \emph{menor} que o menor deles.}
\end{questao}

\begin{questao}[1]{}
Dois capacitores de $\SI{4}{\micro\farad}$ e $\SI{12}{\micro\farad}$ são ligados
em \textbf{paralelo}. A capacitância equivalente é:
\begin{alternativas}[3]
\item \SI{3}{\micro\farad}
\item \SI{8}{\micro\farad}
\item \SI{16}{\micro\farad}
\item \SI{48}{\micro\farad}
\item \SI{6}{\micro\farad}
\end{alternativas}
\resposta{(c)}
\resolucao{No paralelo, as capacitâncias \emph{somam}:
$\Ceq = 4 + 12 = \SI{16}{\micro\farad}$.}
\end{questao}

\begin{questao}[1]{}
Um capacitor de $\SI{100}{\micro\farad}$ é carregado a $\SI{50}{\volt}$. A energia
armazenada é:
\begin{alternativas}[3]
\item \SI{2.5}{\milli\joule}
\item \SI{5}{\milli\joule}
\item \SI{125}{\milli\joule}
\item \SI{250}{\milli\joule}
\item \SI{5}{\joule}
\end{alternativas}
\resposta{(c)}
\resolucao{$E = \tfrac12 C U^2 = \tfrac12\cdot\pot{100}{-6}\cdot50^2
= \tfrac12\cdot\pot{100}{-6}\cdot2500 = \SI{0.125}{\joule} = \SI{125}{\milli\joule}$.
Essa energia fica armazenada no campo elétrico entre as placas e pode ser liberada
rapidamente --- princípio do flash fotográfico e do desfibrilador.}
\end{questao}

\blocoaplicacao

\begin{questao}[2]{}
Uma associação mista tem $C_1 = \SI{6}{\micro\farad}$ e
$C_2 = \SI{3}{\micro\farad}$ em \emph{paralelo}, e esse conjunto em \emph{série}
com $C_3 = \SI{2}{\micro\farad}$. Calcule a capacitância total.

\circuitoq{%
\begin{circuitikz}[scale=1,transform shape,line width=0.8pt]
 \draw (0,0.75) to[short,o-] (1,0.75) coordinate(a);
 \draw (a) -- (1,1.5) to[C,l={$C_1=\SI{6}{\micro\farad}$}] (3,1.5) -- (3,0.75) coordinate(b);
 \draw (a) -- (1,0) to[C,l_={$C_2=\SI{3}{\micro\farad}$}] (3,0) -- (3,0.75);
 \draw (b) to[C,l={$C_3=\SI{2}{\micro\farad}$}] (5,0.75) to[short,-o] (6,0.75);
\end{circuitikz}}

\resposta{$\Ceq \approx \SI{1.64}{\micro\farad}$}
\resolucao{\textbf{Paralelo} $C_1$ e $C_2$: somam-se,
$C_{12} = 6 + 3 = \SI{9}{\micro\farad}$.\\
\textbf{Série} com $C_3$:
\[
\Ceq=\frac{C_{12}\,C_3}{C_{12}+C_3}=\frac{9\cdot2}{11}=\frac{18}{11}
\approx\SI{1.64}{\micro\farad}.
\]}
\end{questao}

\begin{questao}[2]{}
Uma associação é feita com dois capacitores de $\SI{8}{\micro\farad}$ em
\emph{série}, e esse conjunto em \emph{paralelo} com um capacitor de
$\SI{4}{\micro\farad}$. Calcule a capacitância equivalente.
\resposta{$\Ceq = \SI{8}{\micro\farad}$}
\resolucao{\textbf{Série} dos dois de \SI{8}{\micro\farad}:
$\dfrac{8}{2} = \SI{4}{\micro\farad}$.\\
\textbf{Paralelo} com o de \SI{4}{\micro\farad}:
$4 + 4 = \SI{8}{\micro\farad}$.}
\end{questao}

\begin{questao}[2]{}
Um circuito RC série tem $R = \SI{1}{\kilo\ohm}$ e $C = \SI{1}{\micro\farad}$,
ligado a uma fonte de $\SI{10}{\volt}$ em $t = 0$.
\begin{itensq}
\item Qual a constante de tempo?
\item Qual a tensão no capacitor após um intervalo igual a $\tau$?
\item Qual a tensão em regime permanente?
\end{itensq}
\resposta{(a) $\tau = \SI{1}{\milli\second}$; (b) $\approx\SI{6.32}{\volt}$;
(c) \SI{10}{\volt}}
\resolucao{(a) $\tau = RC = \pot{1}{3}\cdot\pot{1}{-6} = \SI{1}{\milli\second}$.\\
(b) Na carga, $v_C(\tau) = U_0(1 - e^{-1}) = 10\cdot0{,}632 \approx \SI{6.32}{\volt}$
--- o capacitor atinge cerca de $\SI{63}{\percent}$ da tensão final em um $\tau$.\\
(c) Em regime permanente ($t\to\infty$), o capacitor está totalmente carregado,
não circula corrente, e $v_C = U_0 = \SI{10}{\volt}$: o capacitor se comporta como
um circuito \emph{aberto}.}
\end{questao}

\begin{questao}[2]{}
No circuito da figura, determine a tensão no capacitor em regime permanente.

\circuitoq{%
\begin{circuitikz}[scale=1,transform shape,line width=0.8pt]
 \draw (0,0) to[battery1,l_={$\SI{12}{\volt}$}] (0,3)
       to[R,l={$R_1=\SI{4}{\ohm}$}] (3,3) coordinate(a);
 \draw (a) to[R,l={$R_2=\SI{2}{\ohm}$}] (3,0) coordinate(b) -- (0,0);
 \draw (a) -- (5,3) to[C,l={$C$}] (5,0) -- (b);
 \node[right,Vcol,font=\footnotesize] at (5.1,1.5) {$v_C$};
\end{circuitikz}}

\resposta{$v_C = \SI{4}{\volt}$}
\resolucao{Em regime permanente, o capacitor está carregado e \emph{não conduz}
(ramo aberto). Toda a corrente passa por $R_1$ e $R_2$ em série:
\[
I=\frac{12}{4+2}=\SI{2}{\ampere}.
\]
A tensão no capacitor iguala a tensão sobre $R_2$ (estão em paralelo):
\[
v_C = R_2\,I = 2\cdot2 = \SI{4}{\volt}.
\]
\textbf{Princípio-chave:} em regime permanente CC, capacitor = circuito aberto
(e indutor = curto-circuito). Isso transforma o problema em uma simples análise
resistiva.}
\end{questao}

\begin{questao}[2]{}
Três capacitores de $\SI{6}{\micro\farad}$ cada são associados. Determine a
capacitância equivalente se estiverem ligados (a) todos em série; (b) todos em
paralelo; (c) dois em paralelo, e o conjunto em série com o terceiro.
\resposta{(a) \SI{2}{\micro\farad}; (b) \SI{18}{\micro\farad}; (c) \SI{4}{\micro\farad}}
\resolucao{(a) Série de três iguais: $\Ceq = \dfrac{C}{3} = \SI{2}{\micro\farad}$.\\
(b) Paralelo de três iguais: $\Ceq = 3C = \SI{18}{\micro\farad}$.\\
(c) Dois em paralelo: $6+6 = \SI{12}{\micro\farad}$; em série com o terceiro:
$\dfrac{12\cdot6}{18} = \SI{4}{\micro\farad}$.\\
Repare como a topologia muda drasticamente o resultado, de \SI{2}{} a
\SI{18}{\micro\farad} com os mesmos três componentes.}
\end{questao}

\blocoaprofundamento

\begin{questao}[3]{}
Um capacitor de $\SI{2}{\micro\farad}$ está carregado a $\SI{6}{\volt}$ e é
descarregado através de um resistor de $\SI{500}{\ohm}$.
\begin{itensq}
\item Qual a constante de tempo da descarga?
\item Qual a corrente inicial de descarga?
\item Qual a carga restante no capacitor após um intervalo igual a $\tau$?
\end{itensq}
\resposta{(a) $\tau = \SI{1}{\milli\second}$; (b) $I_0 = \SI{12}{\milli\ampere}$;
(c) $\approx\SI{4.4}{\micro\coulomb}$}
\resolucao{(a) $\tau = RC = 500\cdot\pot{2}{-6} = \SI{1}{\milli\second}$.\\
(b) A corrente inicial é limitada apenas pelo resistor:
$I_0 = \dfrac{U_0}{R} = \dfrac{6}{500} = \SI{12}{\milli\ampere}$.\\
(c) A carga inicial é $Q_0 = C\,U_0 = \pot{2}{-6}\cdot6 = \SI{12}{\micro\coulomb}$.
Após um $\tau$, resta $Q(\tau) = Q_0\,e^{-1} = 12\cdot0{,}368
\approx\SI{4.4}{\micro\coulomb}$ --- cerca de $\SI{37}{\percent}$ do inicial.}
\end{questao}

\begin{questao}[3]{}
O gráfico mostra a tensão $v_C(t)$ sobre um capacitor durante a carga em um
circuito RC alimentado por $\SI{10}{\volt}$.

\circuitoqq{%
\begin{tikzpicture}
 \begin{axis}[width=9.5cm,height=5cm,xlabel={$t$ (ms)},ylabel={$v_C$ (V)},
   xmin=0,xmax=5,ymin=0,ymax=11,xtick={0,1,2,3,4,5},ytick={0,6.32,10},
   yticklabels={0,6.32,10},grid=both,axis lines=left,domain=0:5,samples=100]
   \addplot[Ccol,very thick] {10*(1-exp(-x))};
   \addplot[dashed,cinza] coordinates {(1,0)(1,6.32)};
   \addplot[dashed,cinza] coordinates {(0,6.32)(1,6.32)};
   \node[Ccol,anchor=west,font=\footnotesize] at (axis cs:1.1,6.0) {$(\tau,\,0{,}63\,U_0)$};
 \end{axis}
\end{tikzpicture}}{Carga de um capacitor em circuito RC.}

Sabendo que a tensão atinge $\SI{6.32}{\volt}$ em $t = \SI{1}{\milli\second}$,
determine a constante de tempo e a tensão final.
\resposta{$\tau = \SI{1}{\milli\second}$; $U_0 = \SI{10}{\volt}$}
\resolucao{O valor $\SI{6.32}{\volt}$ corresponde a $\SI{63}{\percent}$ da tensão
final --- por definição, esse é o instante $t = \tau$. Logo
$\tau = \SI{1}{\milli\second}$. A curva satisfaz $v_C = U_0(1 - e^{-t/\tau})$, e
sua assíntota horizontal é a tensão da fonte, $U_0 = \SI{10}{\volt}$. Lendo o
gráfico: a $\SI{63}{\percent}$ em $\SI{1}{\milli\second}$ e saturando em
\SI{10}{\volt}, confirma-se $\tau = \SI{1}{\milli\second}$.}
\end{questao}

\begin{questao}[3]{}
Um capacitor $C_1 = \SI{10}{\micro\farad}$ é carregado a $\SI{100}{\volt}$ e
depois conectado, por meio de uma chave, a um capacitor descarregado
$C_2 = \SI{5}{\micro\farad}$.

\circuitoq{%
\begin{circuitikz}[scale=1,transform shape,line width=0.8pt]
 \draw (0,0) to[C,l_={$C_1$}] (0,2.6) -- (1.6,2.6) coordinate(a);
 \draw (a) to[cspst,l={S}] (3.4,2.6) -- (4.6,2.6);
 \draw (4.6,2.6) to[C,l={$C_2$}] (4.6,0) -- (0,0);
 \node[font=\footnotesize] at (0,-0.45) {carregado a \SI{100}{\volt}};
 \node[font=\footnotesize] at (4.6,-0.45) {descarregado};
\end{circuitikz}}

Após fechar a chave:
\begin{itensq}
\item Determine a tensão final comum aos dois capacitores.
\item Calcule a energia armazenada antes e depois.
\item Explique o "paradoxo": para onde foi a energia que desapareceu?
\end{itensq}
\resposta{(a) $V_f \approx \SI{66.7}{\volt}$; (b) \SI{50}{\milli\joule} $\to$
\SI{33.3}{\milli\joule}; (c) ver comentário}
\resolucao{\textbf{(a)} A carga total se \emph{conserva} e se redistribui até que
as tensões se igualem:
\[
Q_0 = C_1V_0=(\pot{10}{-6})(100)=\SI{1}{\milli\coulomb},
\]
\[
V_f=\frac{Q_0}{C_1+C_2}=\frac{\pot{1}{-3}}{\pot{15}{-6}}
=\frac{200}{3}\approx\SI{66.7}{\volt}.
\]
\textbf{(b) Energias:}
\[
E_i=\tfrac12 C_1V_0^{\,2}=\tfrac12(\pot{10}{-6})(100)^2=\SI{50}{\milli\joule},
\]
\[
E_f=\tfrac12 (C_1+C_2)V_f^{\,2}
=\tfrac12(\pot{15}{-6})\left(\tfrac{200}{3}\right)^{2}
=\frac{100}{3}\approx\SI{33.3}{\milli\joule}.
\]
Perderam-se $\SI{16.7}{\milli\joule}$ --- exatamente $\dfrac{1}{3}$ da energia
inicial.

\textbf{(c) Resolução do paradoxo.} A carga se conserva, mas a energia
\textbf{não} precisa se conservar \emph{dentro do subsistema dos capacitores} ---
ela foi \emph{dissipada}. Os mecanismos:
\begin{itemize}[leftmargin=1.6em,itemsep=1pt,topsep=2pt]
\item \textbf{Resistência dos fios e da chave:} por menor que seja, a corrente
      transitória a atravessa e gera calor $\int Ri^2\,\mathrm{d}t$. O
      resultado notável é que \emph{a energia perdida independe do valor de $R$}
      --- só o \emph{tempo} do transitório muda.
\item \textbf{Radiação eletromagnética:} no limite $R\to0$, o circuito LC
      oscilaria e irradiaria a energia como ondas.
\end{itemize}
\textbf{Fórmula geral da perda:}
\[
\Delta E=\frac{1}{2}\,\frac{C_1C_2}{C_1+C_2}\,V_0^{\,2}
=\frac12\cdot\frac{10\cdot5}{15}\cdot10^{-6}\cdot 100^2
\approx\SI{16.7}{\milli\joule}.\ \checkmark
\]
\textbf{Relevância prática:} esse é o motivo pelo qual bancos de capacitores nunca
são conectados em paralelo sem resistores de pré-carga (\emph{soft-start}) ---
a corrente de equalização pode ser destrutiva.}
\end{questao}

\begin{questao}[3]{}
Um capacitor de placas paralelas com $C_0 = \SI{10}{\micro\farad}$ é carregado a
$U_0 = \SI{100}{\volt}$ e então \textbf{desconectado} da fonte. Em seguida,
introduz-se entre as placas um dielétrico de constante $\kappa = 4$.
\begin{itensq}
\item O que acontece com a carga, a capacitância, a tensão e a energia?
\item Calcule os valores finais de cada grandeza.
\item Para onde foi a energia "perdida"? Ela violou a conservação de energia?
\item E se o dielétrico fosse inserido com o capacitor \emph{ainda ligado} à
      fonte? Compare os dois casos.
\end{itensq}
\resposta{(b) $Q=\SI{1}{\milli\coulomb}$ (constante), $C=\SI{40}{\micro\farad}$,
$U=\SI{25}{\volt}$, $E$ cai de \SI{50}{\milli\joule} para \SI{12.5}{\milli\joule};
(c) e (d) ver comentário}
\resolucao{\textbf{(a) e (b) --- capacitor isolado ($Q$ constante).}
Desconectado, a carga não tem para onde ir:
\[
Q = C_0U_0 = \pot{10}{-6}\cdot100 = \SI{1}{\milli\coulomb}\ \text{(constante)}.
\]
O dielétrico multiplica a capacitância por $\kappa$:
\[
C = \kappa C_0 = \SI{40}{\micro\farad},
\qquad
U = \frac{Q}{C}=\frac{\pot{1}{-3}}{\pot{40}{-6}}=\SI{25}{\volt}
\quad\left(=\frac{U_0}{\kappa}\right).
\]
Energias:
\[
E_i=\tfrac12 C_0U_0^{\,2}=\SI{50}{\milli\joule},
\qquad
E_f=\tfrac12 CU^{2}=\tfrac12(\pot{40}{-6})(25)^2=\SI{12.5}{\milli\joule}
=\frac{E_i}{\kappa}.
\]
\textbf{(c) Para onde foi a energia?} \emph{Não houve violação alguma.} O
dielétrico é \textbf{atraído} para dentro do capacitor pelo campo de borda ---
existe uma força elétrica real puxando-o. Se ele for solto, essa força realiza
trabalho e o acelera; a energia "perdida" pelo campo converte-se em energia
cinética do dielétrico (e, ao final, em calor por atrito e vibração).
Para inseri-lo \emph{quase-estaticamente}, um agente externo precisa segurá-lo,
e é ele quem \emph{recebe} os \SI{37.5}{\milli\joule}.

\textbf{(d) Caso alternativo --- capacitor ligado à fonte ($U$ constante).}
Agora a tensão é fixada em \SI{100}{\volt}:
\[
C = \SI{40}{\micro\farad},\qquad
Q = CU = \SI{4}{\milli\coulomb}\ (\text{quadruplicou}),\qquad
E_f=\tfrac12 CU^{2}=\SI{200}{\milli\joule}\ (\text{quadruplicou}).
\]
A energia armazenada \emph{aumentou} \SI{150}{\milli\joule}. A fonte, porém,
forneceu $\Delta Q\cdot U = (\pot{3}{-3})(100)=\SI{300}{\milli\joule}$ --- o
dobro. A metade restante (\SI{150}{\milli\joule}) foi entregue como trabalho
mecânico ao dielétrico (que também é sugado para dentro).

\smallskip
\begin{center}\small\sffamily
\begin{tabular}{@{}l c c@{}}
\toprule
 & \textbf{Isolado ($Q$ const.)} & \textbf{Ligado ($U$ const.)}\\
\midrule
Carga & constante & $\times\kappa$\\
Tensão & $\div\kappa$ & constante\\
Energia & $\div\kappa$ (diminui) & $\times\kappa$ (aumenta)\\
\bottomrule
\end{tabular}
\end{center}
\textbf{Moral:} a mesma ação física (inserir o dielétrico) produz resultados
\emph{opostos} para a energia conforme a condição de contorno. Identificar
\emph{qual grandeza permanece constante} é o passo decisivo do problema.}
\end{questao}

% BEGIN BANCO COMPLEMENTAR Capacitores
% =====================================================================
%  Banco complementar --- Capacitores  (REVISADO)
%  Substitui a versao gerada por template (10 clones em N2, 11 em N3,
%  10 questoes conceituais indevidamente marcadas como N4).
%  Sem sobreposicao com o cap10, que ja traz: C de placas paralelas,
%  Q=CV com W, serie 6/3, paralelo 4/12, energia (MC), associacao mista
%  6||3, RC de 1 kohm/1 uF, tres de 6 uF, descarga de 2 uF/6 V, grafico
%  v_C(t), COMPARTILHAMENTO DE CARGA entre C1 e C2 e INSERCAO DE
%  DIELETRICO (a Q e a V constantes) --- estes dois ultimos motivaram a
%  remocao de duas N4 do banco antigo, que os duplicavam.
%  Distribuicao mantida: 4 N1 + 10 N2 + 11 N3 + 10 N4 = 35
% =====================================================================

\subsubsection*{Banco complementar --- Capacitores}

\begin{atencao}[Organização do banco]
Três relações organizam o tópico: $Q=CU$, $W=\frac12CU^{2}$ e
$i=C\dfrac{\mathrm{d}u}{\mathrm{d}t}$. Delas decorre o fato central dos
transitórios --- a \textbf{tensão} no capacitor não pode saltar (isso exigiria
corrente infinita), ao contrário do indutor, em que é a \emph{corrente} que é
contínua. O N3 trata de divisão de tensão em série, fuga, ESR e extração de
parâmetros; o N4 exige dedução da EDO, balanço energético, projeto de banco,
filtro e temporizador.
\end{atencao}

\blocofixacao

\begin{questao}[1]{}
Um capacitor armazena $\SI{300}{\micro\coulomb}$ sob $\SI{20}{\volt}$.
Determine sua capacitância.
\resposta{$C=\SI{15}{\micro\farad}$}
\resolucao{
\[
C=\frac{Q}{U}=\frac{\pot{300}{-6}}{20}=\pot{15}{-6}\,\si{\farad}
=\SI{15}{\micro\farad}.
\]
A capacitância é uma propriedade \emph{geométrica} do componente: ela não
depende de $Q$ nem de $U$, apenas da razão entre eles. Aplicar o dobro da
tensão dobra a carga e deixa $C$ inalterada.}
\end{questao}

\begin{questao}[1]{}
Um capacitor de $\SI{100}{\micro\farad}$ armazena $\SI{8}{\milli\joule}$.
Determine a tensão em seus terminais.
\resposta{$U\approx\SI{12.6}{\volt}$}
\resolucao{De $W=\frac12CU^{2}$,
\[
U=\sqrt{\frac{2W}{C}}=\sqrt{\frac{2(\pot{8}{-3})}{\pot{100}{-6}}}
=\sqrt{160}=\SI{12.65}{\volt}.
\]
A raiz quadrada indica que a energia é pouco sensível à tensão em valores
baixos e muito sensível em valores altos: dobrar a tensão \emph{quadruplica}
a energia armazenada.}
\end{questao}

\begin{questao}[1]{}
Sobre a tensão em um capacitor ideal, é correto afirmar que ela:
\begin{alternativas}[2]
\item pode variar instantaneamente
\item não pode variar instantaneamente
\item é sempre igual à da fonte
\item é sempre nula em regime permanente CC
\end{alternativas}
\resposta{(b)}
\resolucao{Um salto de tensão implicaria
$\mathrm{d}u/\mathrm{d}t\to\infty$ e, por
$i=C\,\mathrm{d}u/\mathrm{d}t$, corrente infinita --- impossível. Logo
$u(0^{+})=u(0^{-})$.\\
\textbf{Dualidade com o indutor:} no capacitor é a \emph{tensão} que é
contínua; no indutor, a \emph{corrente}. Em regime CC o capacitor tem corrente
nula (comporta-se como circuito aberto) e tensão em geral não nula --- o que
descarta a alternativa (d).}
\end{questao}

\begin{questao}[1]{}
Determine a constante de tempo de um circuito RC com
$R=\SI{4.7}{\kilo\ohm}$ e $C=\SI{10}{\micro\farad}$.
\resposta{$\tau=\SI{47}{\milli\second}$}
\resolucao{
\[
\tau=RC=4700\cdot\pot{10}{-6}=\pot{4.7}{-2}\,\si{\second}
=\SI{47}{\milli\second}.
\]
\textbf{Atenção à estrutura:} no RC, resistência \emph{maior} torna o
transitório \emph{mais lento} ($\tau=RC$); no RL ocorre o contrário
($\tau=L/R$). Verificar a dimensão ajuda:
$\si{\ohm}\times\si{\farad}=\si{\second}$.}
\end{questao}

\blocoaplicacao

\begin{questao}[2]{}
Dois capacitores de $\SI{10}{\micro\farad}$ em série são associados em paralelo
com um de $\SI{15}{\micro\farad}$. Determine $C_{eq}$.
\resposta{$C_{eq}=\SI{20}{\micro\farad}$}
\resolucao{\textbf{Série} (regra inversa, ao contrário dos resistores):
\[
C_s=\frac{10\cdot10}{20}=\SI{5}{\micro\farad}.
\]
\textbf{Paralelo} (soma direta):
\[
C_{eq}=5+15=\SI{20}{\micro\farad}.
\]
Memorize a inversão: capacitores em \emph{série} combinam como resistores em
paralelo, e vice-versa. A razão física é que a série reduz a capacitância
(equivale a afastar as placas), enquanto o paralelo a aumenta (equivale a somar
áreas).}
\end{questao}

\begin{questao}[2]{}
Capacitores de $\SI{1}{\micro\farad}$ e $\SI{2}{\micro\farad}$ estão em série
sob $\SI{30}{\volt}$. Determine $C_{eq}$, a carga e a tensão em cada um.
\resposta{$C_{eq}=\SI{0.667}{\micro\farad}$; $Q=\SI{20}{\micro\coulomb}$;
$U_1=\SI{20}{\volt}$, $U_2=\SI{10}{\volt}$}
\resolucao{
\[
C_{eq}=\frac{1\cdot2}{3}=\SI{0.667}{\micro\farad};
\qquad
Q=C_{eq}U=0{,}667\cdot30=\SI{20}{\micro\coulomb}.
\]
Em série, a \textbf{carga é a mesma} em todos (é a mesma corrente que os
carregou):
\[
U_1=\frac{20}{1}=\SI{20}{\volt},
\qquad
U_2=\frac{20}{2}=\SI{10}{\volt}.
\]
\textbf{Verificação:} $20+10=\SI{30}{\volt}$ \checkmark.\\
A tensão se reparte \emph{inversamente} à capacitância --- o capacitor
\textbf{menor} fica com a maior tensão. É o oposto da série de resistores, e é
a causa de um problema de projeto tratado adiante.}
\end{questao}

\begin{questao}[2]{}
A tensão sobre um capacitor de $\SI{10}{\micro\farad}$ cresce linearmente à
taxa de $\SI{500}{\volt\per\second}$. Determine a corrente.
\resposta{$i=\SI{5}{\milli\ampere}$}
\resolucao{
\[
i=C\frac{\mathrm{d}u}{\mathrm{d}t}=\pot{10}{-6}\cdot500
=\pot{5}{-3}\,\si{\ampere}=\SI{5}{\milli\ampere}.
\]
Enquanto a rampa durar, a corrente é \emph{constante}. Se a tensão parasse de
variar, a corrente cairia a zero instantaneamente --- o capacitor responde à
\textbf{variação} da tensão, não ao seu valor.}
\end{questao}

\begin{questao}[2]{}
Aplica-se a um capacitor de $\SI{2.2}{\micro\farad}$ uma tensão
\textbf{triangular} que sobe $\SI{10}{\volt}$ em $\SI{1}{\milli\second}$ e
desce igualmente no milissegundo seguinte. Descreva e quantifique a corrente.
\resposta{Onda \textbf{quadrada} de $\pm\SI{22}{\milli\ampere}$}
\resolucao{Na subida,
$\mathrm{d}u/\mathrm{d}t=10/10^{-3}=\SI{10000}{\volt\per\second}$:
\[
i=2{,}2\times10^{-6}\cdot10^{4}=\SI{22}{\milli\ampere}.
\]
Na descida a taxa inverte e $i=\SI{-22}{\milli\ampere}$.\\
\textbf{Resultado geral:} o capacitor \emph{deriva} a tensão --- triangular
entra, quadrada sai. O indutor faz o oposto (integra): tensão quadrada nele
produz corrente triangular. É a dualidade dos dois elementos em sua forma mais
visível no osciloscópio.}
\end{questao}

\begin{questao}[2]{}
Um circuito RC com $\tau=\SI{2}{\milli\second}$ é ligado a
$\SI{20}{\volt}$ com o capacitor inicialmente descarregado. Determine
$u_C$ em $t=\tau$, $2\tau$ e $3\tau$.
\resposta{$\SI{12.64}{\volt}$, $\SI{17.29}{\volt}$ e $\SI{19.00}{\volt}$}
\resolucao{Com $u_C=U\left(1-e^{-t/\tau}\right)$ e $U=\SI{20}{\volt}$:
\[
u(\tau)=20(0{,}632)=\SI{12.64}{\volt};
\quad
u(2\tau)=20(0{,}865)=\SI{17.29}{\volt};
\quad
u(3\tau)=20(0{,}950)=\SI{19.00}{\volt}.
\]
As frações $63{,}2\%$, $86{,}5\%$ e $95{,}0\%$ são as mesmas do circuito RL ---
elas dependem só da forma exponencial, não do tipo de elemento.}
\end{questao}

\begin{questao}[2]{}
No mesmo circuito, determine o instante em que a tensão atinge
$\SI{15}{\volt}$.
\resposta{$t\approx\SI{2.77}{\milli\second}$}
\resolucao{
\[
15=20\left(1-e^{-t/\tau}\right)
\;\Longrightarrow\;
e^{-t/\tau}=0{,}25
\;\Longrightarrow\;
t=\tau\ln4=2(1{,}386)=\SI{2.77}{\milli\second}.
\]
\textbf{Atalho útil:} atingir $75\%$ leva exatamente $\ln4\approx1{,}39\tau$;
atingir $50\%$ leva $\ln2\approx0{,}69\tau$. Esse último valor é a "meia-vida"
do transitório e vale a pena memorizar.}
\end{questao}

\begin{questao}[2]{}
Uma fonte de $\SI{24}{\volt}$ alimenta $R_1=\SI{8}{\kilo\ohm}$ em série com
$R_2=\SI{4}{\kilo\ohm}$. Um capacitor está ligado em paralelo com $R_2$.
Determine a tensão no capacitor em regime permanente e a corrente que o
atravessa.
\resposta{$U_C=\SI{8}{\volt}$; $i_C=0$}
\resolucao{Em regime CC o capacitor é um \textbf{circuito aberto}: não conduz.
Toda a corrente passa pelos dois resistores, e a tensão do capacitor é a de
$R_2$:
\[
U_C=24\cdot\frac{4}{8+4}=\SI{8}{\volt};
\qquad
i_C=0.
\]
\textbf{Método:} para determinar o estado final de qualquer circuito CC com
capacitores, basta \emph{removê-los} do esquema e resolver a rede resistiva
restante. Para indutores, o procedimento dual é substituí-los por curtos.}
\end{questao}

\begin{questao}[2]{}
Dois capacitores armazenam a mesma energia. O primeiro tem
$C_1=\SI{100}{\micro\farad}$ sob $\SI{10}{\volt}$. Se
$C_2=\SI{25}{\micro\farad}$, determine sua tensão.
\resposta{$U_2=\SI{20}{\volt}$}
\resolucao{
\[
\frac12C_1U_1^{2}=\frac12C_2U_2^{2}
\;\Longrightarrow\;
U_2=U_1\sqrt{\frac{C_1}{C_2}}=10\sqrt{4}=\SI{20}{\volt}.
\]
Ambos armazenam $\SI{5}{\milli\joule}$. Reduzir $C$ por um fator $4$ exige
\emph{dobrar} a tensão --- e a tensão é justamente o parâmetro caro e limitado
em capacitores reais, pois exige dielétrico mais espesso, o que reduz $C$.}
\end{questao}

\begin{questao}[2]{}
Um capacitor plano tem capacitância $C$. Determine a nova capacitância se:
\begin{itensq}
\item a área das placas for dobrada;
\item a distância entre placas for reduzida à metade;
\item ambas as alterações forem feitas.
\end{itensq}
\resposta{(a) $2C$; (b) $2C$; (c) $4C$}
\resolucao{De $C=\dfrac{\varepsilon A}{d}$: $C\propto A$ e
$C\propto1/d$.\\
(a) $A\to2A$ dá $2C$. (b) $d\to d/2$ dá $2C$. (c) os efeitos se multiplicam:
$4C$.\\
\textbf{Limite prático de (b):} reduzir $d$ aumenta o campo
$E=U/d$ para a mesma tensão. Chega-se rapidamente à rigidez dielétrica do
material, e o capacitor perfura. Há, portanto, um teto para essa estratégia ---
motivo pelo qual capacitores de alta capacitância usam grandes áreas
(folhas enroladas) em vez de dielétricos ultrafinos.}
\end{questao}

\begin{questao}[2]{}
Um capacitor carregado a $\SI{50}{\volt}$ é descarregado através de
$R=\SI{2.2}{\kilo\ohm}$. Determine a corrente inicial e explique sua evolução.
\resposta{$i_0=\SI{22.7}{\milli\ampere}$, decaindo exponencialmente}
\resolucao{No instante inicial a tensão ainda é $\SI{50}{\volt}$
(continuidade), e ela se aplica integralmente ao resistor:
\[
i_0=\frac{50}{2200}=\SI{22.7}{\milli\ampere}.
\]
À medida que o capacitor se descarrega, sua tensão cai e a corrente a acompanha:
$i(t)=i_0e^{-t/\tau}$.\\
\textbf{Observação:} a corrente é máxima no início e a tensão é máxima no
início --- diferentemente do circuito RL, em que a corrente parte de zero. Essa
diferença decorre justamente de qual grandeza é contínua em cada elemento.}
\end{questao}

\blocoaprofundamento

\begin{questao}[3]{}
$C_1=\SI{4}{\micro\farad}$ e $C_2=\SI{12}{\micro\farad}$, ambos com tensão
máxima de trabalho de $\SI{50}{\volt}$, são ligados em série.
\begin{itensq}
\item Determine $C_{eq}$ e a distribuição de tensão sob $\SI{100}{\volt}$.
\item Determine a máxima tensão que a associação suporta.
\item Explique por que o resultado é contraintuitivo.
\end{itensq}
\resposta{(a) $\SI{3}{\micro\farad}$; $\SI{75}{\volt}$ e $\SI{25}{\volt}$;
(b) $\SI{66.7}{\volt}$}
\resolucao{(a) $C_{eq}=\dfrac{4\cdot12}{16}=\SI{3}{\micro\farad}$. Como a carga
é comum, $U\propto1/C$:
\[
U_1=100\cdot\frac{12}{16}=\SI{75}{\volt},
\qquad
U_2=100\cdot\frac{4}{16}=\SI{25}{\volt}.
\]
(b) Sob $\SI{100}{\volt}$, $C_1$ já está a $\SI{75}{\volt}$ --- $50\%$ acima do
limite. Impondo $U_1\le\SI{50}{\volt}$:
\[
U_{\max}=50\cdot\frac{16}{12}=\SI{66.7}{\volt}.
\]
(c) \textbf{É contraintuitivo} porque associar dois componentes de
$\SI{50}{\volt}$ em série sugere $\SI{100}{\volt}$ de capacidade --- e isso
\emph{só} é verdade se eles forem \textbf{iguais}. Com valores diferentes, o
menor capacitor absorve a maior parte da tensão e limita todo o conjunto.\\
\textbf{Regra prática:} em série de capacitores, use sempre unidades idênticas
--- e ainda assim adicione resistores de equalização, como se detalha no bloco
de desafio.}
\end{questao}

\begin{questao}[3]{}
Capacitores de $\SI{2}{\micro\farad}$, $\SI{3}{\micro\farad}$ e
$\SI{5}{\micro\farad}$ estão em paralelo sob $\SI{20}{\volt}$. Determine a
carga de cada um, a carga total e $C_{eq}$, e compare com o caso série.
\resposta{$40$, $60$ e $\SI{100}{\micro\coulomb}$; total
$\SI{200}{\micro\coulomb}$; $C_{eq}=\SI{10}{\micro\farad}$}
\resolucao{Em paralelo a \textbf{tensão} é comum:
\[
Q_k=C_kU \Rightarrow 40,\ 60,\ \SI{100}{\micro\coulomb};
\qquad
Q_{tot}=\SI{200}{\micro\coulomb}.
\]
\[
C_{eq}=2+3+5=\SI{10}{\micro\farad}
\quad\text{(confere: }200/20=10\text{)}.
\]
\textbf{Comparação com a série:} lá, a \emph{carga} é comum e a tensão se
divide; aqui, a \emph{tensão} é comum e a carga se divide. O paralelo soma
capacitâncias e mantém a tensão de trabalho da menor unidade; a série reduz a
capacitância e (idealmente) soma as tensões. São estratégias opostas para
problemas opostos.}
\end{questao}

\begin{questao}[3]{}
Uma fonte de $\SI{12}{\volt}$ alimenta $R_1=\SI{6}{\kilo\ohm}$ em série com
$R_2=\SI{3}{\kilo\ohm}$, e um capacitor de $\SI{100}{\nano\farad}$ está em
paralelo com $R_2$.
\begin{itensq}
\item Determine a tensão final no capacitor.
\item Determine $\tau$ pelo equivalente de Thévenin.
\item Escreva $u_C(t)$.
\end{itensq}
\resposta{(a) $\SI{4}{\volt}$; (b) $\tau=\SI{0.2}{\milli\second}$;
(c) $u_C=4\left(1-e^{-t/0{,}2\,\mathrm{ms}}\right)$}
\resolucao{(a) Com o capacitor aberto em regime,
$U_C=12\cdot\dfrac{3}{9}=\SI{4}{\volt}$.\\
(b) Zerando a fonte, o capacitor vê $R_1\parallel R_2$:
\[
R_{Th}=\frac{6\cdot3}{9}=\SI{2}{\kilo\ohm}
\;\Longrightarrow\;
\tau=R_{Th}C=2000\cdot\pot{100}{-9}=\SI{0.2}{\milli\second}.
\]
(c) $u_C(t)=4\left(1-e^{-t/0{,}0002}\right)$ volts.\\
\textbf{Erro comum:} usar $\tau=R_1C$ ou $\tau=R_2C$. A constante de tempo é
governada pela resistência \emph{vista pelo capacitor} com as fontes zeradas
--- aqui, o paralelo das duas, que é menor que qualquer uma delas. O transitório
é, portanto, mais rápido do que qualquer estimativa ingênua sugeriria.}
\end{questao}

\begin{questao}[3]{}
Um capacitor já está a $\SI{2}{\volt}$ quando é ligado a um circuito que o
levaria a $\SI{10}{\volt}$, com $\tau=\SI{5}{\milli\second}$.
\begin{itensq}
\item Escreva $u_C(t)$.
\item Determine a tensão em $t=\SI{5}{\milli\second}$.
\item Determine a corrente inicial, sabendo que $R=\SI{5}{\kilo\ohm}$.
\end{itensq}
\resposta{(a) $u_C=10-8e^{-t/\tau}$; (b) $\SI{7.06}{\volt}$;
(c) $\SI{1.6}{\milli\ampere}$}
\resolucao{(a) A forma geral de qualquer transitório de primeira ordem é
\[
u(t)=U_f+\left(U_0-U_f\right)e^{-t/\tau}=10-8e^{-t/\tau}.
\]
(b) $u(5\,\mathrm{ms})=10-8e^{-1}=10-2{,}94=\SI{7.06}{\volt}$.\\
(c) Em $t=0^{+}$ a tensão do capacitor é $\SI{2}{\volt}$, e sobre o resistor
sobram $10-2=\SI{8}{\volt}$:
\[
i_0=\frac{8}{5000}=\SI{1.6}{\milli\ampere}.
\]
\textbf{Vantagem da forma geral:} ela cobre carga ($U_0=0$), descarga
($U_f=0$) e qualquer condição intermediária. Basta identificar os três
parâmetros $U_0$, $U_f$ e $\tau$ --- não há três fórmulas a memorizar.}
\end{questao}

\begin{questao}[3]{}
Um capacitor de $\SI{470}{\micro\farad}$ tem resistência de isolação de
$\SI{2}{\mega\ohm}$. Carregado a $\SI{100}{\volt}$ e deixado em circuito
aberto:
\begin{itensq}
\item Determine a constante de tempo de autodescarga.
\item Determine a tensão após $\SI{15}{\minute}$.
\item Comente a implicação de segurança.
\end{itensq}
\resposta{(a) $\tau=\SI{940}{\second}\approx\SI{15.7}{\minute}$;
(b) $\approx\SI{38}{\volt}$}
\resolucao{(a) A fuga é modelada por um resistor em paralelo:
\[
\tau=R_{iso}C=\pot{2}{6}\cdot\pot{470}{-6}=\SI{940}{\second}.
\]
(b) Com $t=\SI{900}{\second}$:
\[
u=100\,e^{-900/940}=100\,e^{-0{,}957}=\SI{38.4}{\volt}.
\]
(c) \textbf{Depois de quinze minutos desligado, o capacitor ainda guarda
$\SI{38}{\volt}$} --- e, em bancos de fonte chaveada com centenas de volts, o
resíduo é letal por horas. É por isso que equipamentos com barramento CC trazem
\emph{resistores de descarga} (bleeders) permanentes e o aviso de aguardar
antes de abrir o gabinete.\\
\textbf{Nota:} capacitores eletrolíticos ainda apresentam \emph{absorção
dielétrica} --- após descarga completa, a tensão "ressurge" parcialmente em
minutos. Curto-circuitar antes de manusear é procedimento obrigatório, não
precaução exagerada.}
\end{questao}

\begin{questao}[3]{}
Um capacitor eletrolítico tem ESR (resistência equivalente em série) de
$\SI{0.1}{\ohm}$ e conduz corrente de ondulação de
$\SI{1.5}{\ampere}$ eficazes.
\begin{itensq}
\item Determine a potência dissipada internamente.
\item Explique por que essa dissipação limita a vida do componente.
\item Indique duas formas de reduzi-la.
\end{itensq}
\resposta{(a) $P=\SI{0.225}{\watt}$}
\resolucao{(a)
\[
P=\mathrm{ESR}\cdot I_{rms}^{2}=0{,}1(1{,}5)^{2}=\SI{0.225}{\watt}.
\]
(b) Essa potência é dissipada \emph{dentro} do capacitor, elevando a
temperatura do eletrólito. A vida útil de eletrolíticos segue uma regra
empírica severa: ela \textbf{cai pela metade a cada $\SI{10}{\celsius}$} de
elevação. Um componente de $\SI{5000}{\hour}$ a $\SI{105}{\celsius}$ dura
$\SI{20000}{\hour}$ a $\SI{85}{\celsius}$ --- e o inverso também vale.
Pior: o eletrólito evapora, a ESR \emph{aumenta}, a dissipação cresce e o
processo se realimenta.\\
(c) \textbf{(i) Dividir a corrente:} usar $n$ capacitores em paralelo reduz a
corrente por unidade a $I/n$ e a dissipação individual a $P/n^{2}$.
Com dois, cada um dissipa $\SI{56}{\milli\watt}$ --- um quarto do original.\\
\textbf{(ii) Escolher tecnologia de baixa ESR:} capacitores de polímero sólido
têm ESR uma ordem de grandeza menor, ao custo de capacitância e tensão
menores.\\
\textbf{Consequência de diagnóstico:} em fontes chaveadas, o capacitor de saída
é quase sempre o primeiro componente a falhar --- e a ESR crescente, não a
perda de capacitância, é o sintoma que primeiro aparece.}
\end{questao}

\begin{questao}[3]{}
Um capacitor de $\SI{100}{\micro\farad}$, inicialmente descarregado, é
carregado a $\SI{20}{\volt}$ através de um resistor.
\begin{itensq}
\item Determine a energia fornecida pela fonte.
\item Determine a energia armazenada no capacitor.
\item Determine a energia dissipada no resistor e comente.
\end{itensq}
\resposta{(a) $\SI{40}{\milli\joule}$; (b) $\SI{20}{\milli\joule}$;
(c) $\SI{20}{\milli\joule}$ --- exatamente metade}
\resolucao{(a) A fonte mantém $\SI{20}{\volt}$ enquanto entrega toda a carga
final $Q=CU=\SI{2}{\milli\coulomb}$:
\[
W_{fonte}=QU=CU^{2}=\pot{100}{-6}(400)=\SI{40}{\milli\joule}.
\]
(b) $W_C=\frac12CU^{2}=\SI{20}{\milli\joule}$.\\
(c) Por conservação, $W_R=40-20=\SI{20}{\milli\joule}$.\\
\textbf{Resultado notável:} o rendimento do carregamento de um capacitor por
fonte de tensão através de resistor é \textbf{sempre $50\%$} --- não importa o
valor de $R$, nem de $C$, nem da tensão. Diminuir $R$ apenas acelera o
processo; a energia perdida é a mesma. A demonstração formal por integração
está no bloco de desafio.\\
\textbf{Consequência prática:} carregadores eficientes não usam resistor em série
--- usam conversores chaveados, que transferem energia em pacotes através de um
indutor e escapam desse limite de $50\%$.}
\end{questao}

\begin{questao}[3]{}
Um capacitor plano usa dielétrico com rigidez de
$\SI{20}{\kilo\volt\per\milli\meter}$ e espessura de
$\SI{0.1}{\milli\meter}$.
\begin{itensq}
\item Determine a tensão de ruptura.
\item Especifique a tensão nominal com margem de $50\%$.
\item Explique por que a rigidez limita o produto capacitância-tensão.
\end{itensq}
\resposta{(a) $\SI{2}{\kilo\volt}$; (b) $\SI{1}{\kilo\volt}$}
\resolucao{(a) $U_{rup}=E_{rig}\cdot d=20\cdot0{,}1=\SI{2}{\kilo\volt}$.\\
(b) Com margem de $50\%$, especifica-se $\SI{1}{\kilo\volt}$ --- e ainda assim
é prudente derrateá-la em alta temperatura e sob ondulação.\\
(c) Combinando $C=\dfrac{\varepsilon A}{d}$ com
$U_{\max}=E_{rig}d$:
\[
C\,U_{\max}=\varepsilon A E_{rig}
\quad\Longrightarrow\quad
W_{\max}=\frac12CU_{\max}^{2}=\frac12\varepsilon E_{rig}^{2}\,(A\,d).
\]
O termo $Ad$ é o \textbf{volume} de dielétrico. Ou seja, a energia máxima
armazenável depende apenas do volume e das propriedades do material
($\varepsilon$ e $E_{rig}$) --- \emph{não} da geometria escolhida.\\
\textbf{Consequência:} não existe truque geométrico para ganhar energia. Só
avança quem melhora o material. É por isso que o desenvolvimento de capacitores
é, essencialmente, ciência de materiais.}
\end{questao}

\begin{questao}[3]{}
Em um circuito RC de carga, mediu-se tensão final de $\SI{10}{\volt}$,
corrente inicial de $\SI{2}{\milli\ampere}$ e tensão de $\SI{6.32}{\volt}$ em
$\SI{3}{\milli\second}$. Determine $R$ e $C$.
\resposta{$R=\SI{5}{\kilo\ohm}$; $C=\SI{0.6}{\micro\farad}$}
\resolucao{\textbf{Resistência:} em $t=0^{+}$ o capacitor está descarregado e
toda a tensão cai em $R$:
\[
R=\frac{U_f}{i_0}=\frac{10}{\pot{2}{-3}}=\SI{5}{\kilo\ohm}.
\]
\textbf{Constante de tempo:} $6{,}32/10=0{,}632$ --- exatamente a fração de
$1\tau$. Logo $\tau=\SI{3}{\milli\second}$ e
\[
C=\frac{\tau}{R}=\frac{\pot{3}{-3}}{5000}=\pot{6}{-7}\,\si{\farad}
=\SI{0.6}{\micro\farad}.
\]
\textbf{Estrutura da medição:} a corrente inicial isola $R$; o transitório
isola $\tau$; e só a combinação entrega $C$. Nenhuma medida sozinha determina
o capacitor --- é o mesmo padrão do circuito RL, em que o regime permanente dá
$R$ e o transitório dá $L$.}
\end{questao}

\begin{questao}[3]{}
Compare o comportamento de capacitor e indutor preenchendo o raciocínio para os
instantes $t=0^{+}$ e $t\to\infty$, e justifique cada equivalência.
\resposta{Capacitor: curto em $t=0^{+}$ (se descarregado), aberto em regime;
indutor: aberto em $t=0^{+}$ (se sem corrente), curto em regime}
\resolucao{\textbf{Capacitor descarregado em $t=0^{+}$:} $u_C=0$, e um elemento
com tensão nula equivale a um \textbf{curto-circuito}. Em regime CC,
$\mathrm{d}u/\mathrm{d}t=0$ implica $i=0$: equivale a \textbf{circuito
aberto}.\\
\textbf{Indutor sem corrente em $t=0^{+}$:} $i_L=0$, o que equivale a
\textbf{circuito aberto}. Em regime, $\mathrm{d}i/\mathrm{d}t=0$ implica
$v=0$: equivale a \textbf{curto}.\\
\textbf{Atenção à condição inicial:} as equivalências em $t=0^{+}$ pressupõem
estado inicial nulo. Se o capacitor já estiver a $U_0$, ele equivale a uma
\emph{fonte de tensão} $U_0$; se o indutor já conduzir $I_0$, equivale a uma
\emph{fonte de corrente} $I_0$.\\
\textbf{Por que isso é útil:} esses quatro substitutos permitem obter os
valores inicial e final de qualquer transitório de primeira ordem resolvendo
apenas circuitos \emph{resistivos}. Junto com $\tau$, eles determinam a solução
completa sem integrar nada.}
\end{questao}

\begin{questao}[3]{}
Um filtro RC passa-baixas tem $R=\SI{10}{\kilo\ohm}$ e
$C=\SI{47}{\nano\farad}$.
\begin{itensq}
\item Determine a frequência de corte.
\item Relacione-a com a constante de tempo.
\item Descreva o comportamento muito abaixo e muito acima de $f_c$.
\end{itensq}
\resposta{(a) $f_c\approx\SI{339}{\hertz}$; (b) $f_c=1/(2\pi\tau)$}
\resolucao{(a)
\[
f_c=\frac{1}{2\pi RC}=\frac{1}{2\pi(10^{4})(\pot{47}{-9})}
=\SI{339}{\hertz}.
\]
(b) Como $\tau=RC=\SI{0.47}{\milli\second}$, vale
$f_c=\dfrac{1}{2\pi\tau}$: um circuito \emph{lento} (τ grande) filtra a partir
de frequências \emph{baixas}. As duas descrições --- temporal e frequencial ---
são a mesma informação.\\
(c) \textbf{Muito abaixo de $f_c$:} o capacitor tem reatância alta, quase não
carrega, e o sinal passa praticamente intacto (ganho $\approx1$).\\
\textbf{Muito acima:} o capacitor não tem tempo de carregar entre as variações;
a saída é fortemente atenuada, a $-\SI{20}{\decibel}$ por década.\\
\textbf{Exatamente em $f_c$:} o ganho vale $1/\sqrt2\approx0{,}707$
($-\SI{3}{\decibel}$), e a defasagem é de $\SI{45}{\degree}$. É a definição
convencional de frequência de corte.}
\end{questao}

\blocodesafio

\begin{questao}[4]{}
\textbf{(Dedução --- resposta do circuito RC.)} Um RC série com fonte $U$ tem a
chave fechada em $t=0$ e $u_C(0)=0$.
\begin{itensq}
\item Escreva a equação diferencial e resolva por separação de variáveis.
\item Obtenha $i(t)$ e verifique a LKT nos dois extremos.
\item Compare a estrutura da solução com a do circuito RL.
\end{itensq}
\resposta{$u_C=U\left(1-e^{-t/RC}\right)$;
$i=\dfrac{U}{R}e^{-t/RC}$}
\resolucao{(a) Pela LKT, com $i=C\,\mathrm{d}u_C/\mathrm{d}t$:
\[
U=Ri+u_C=RC\frac{\mathrm{d}u_C}{\mathrm{d}t}+u_C
\;\Longrightarrow\;
\frac{\mathrm{d}u_C}{U-u_C}=\frac{\mathrm{d}t}{RC}.
\]
Integrando de $0$ a $t$:
\[
-\ln\!\left(\frac{U-u_C}{U}\right)=\frac{t}{RC}
\;\Longrightarrow\;
u_C(t)=U\left(1-e^{-t/RC}\right).
\]
(b)
\[
i=C\frac{\mathrm{d}u_C}{\mathrm{d}t}
=CU\cdot\frac{1}{RC}e^{-t/RC}=\frac{U}{R}e^{-t/RC}.
\]
\emph{Em $t=0^{+}$:} $u_C=0$ e $Ri=U$; soma $=U$ \checkmark.\\
\emph{Em $t\to\infty$:} $u_C=U$ e $Ri=0$; soma $=U$ \checkmark.\\
(c) \textbf{Dualidade estrutural.} No RL, a \emph{corrente} sobe como
$1-e^{-t/\tau}$ e a tensão do indutor decai como $e^{-t/\tau}$. No RC, é a
\emph{tensão} que sobe e a corrente que decai. Trocando
$L\leftrightarrow C$, $i\leftrightarrow u$ e $R\leftrightarrow1/R$, uma solução
vira a outra --- razão pela qual basta dominar um dos dois casos.}
\end{questao}

\begin{questao}[4]{}
\textbf{(Demonstração --- rendimento do carregamento.)} Um capacitor
descarregado é carregado por uma fonte $U$ através de $R$.
\begin{itensq}
\item Integre $p_R=Ri^{2}$ para obter a energia dissipada.
\item Mostre que ela vale $\frac12CU^{2}$, independentemente de $R$.
\item Explique como conversores chaveados contornam esse limite.
\end{itensq}
\resposta{$W_R=\frac12CU^{2}$; rendimento sempre $50\%$}
\resolucao{(a) Com $i=\dfrac{U}{R}e^{-t/\tau}$ e $\tau=RC$:
\[
W_R=\int_0^{\infty}R\left(\frac{U}{R}\right)^{2}e^{-2t/\tau}\,\mathrm{d}t
=\frac{U^{2}}{R}\left[-\frac{\tau}{2}e^{-2t/\tau}\right]_0^{\infty}
=\frac{U^{2}\tau}{2R}.
\]
(b) Substituindo $\tau=RC$:
\[
W_R=\frac{U^{2}}{2R}\cdot RC=\frac12CU^{2}.
\]
O $R$ \textbf{cancela}: ele aparece na potência ($\propto1/R$) e desaparece na
duração ($\tau\propto R$). Como a fonte forneceu $CU^{2}$ e o capacitor reteve
$\frac12CU^{2}$, o rendimento é exatamente $50\%$ --- sempre.\\
(c) \textbf{O limite decorre de forçar a tensão do capacitor a saltar de $0$
para $U$ contra uma fonte rígida.} Conversores chaveados evitam isso: um
indutor recebe energia da fonte em regime de corrente, e depois a entrega ao
capacitor. Como o indutor é um elemento reativo (não dissipa), a transferência
pode, em princípio, ser feita sem perda --- e na prática rendimentos de $90\%$
a $95\%$ são rotineiros.\\
\textbf{Observação:} o mesmo resultado de $50\%$ aparece no compartilhamento de
carga entre dois capacitores e na descarga de um indutor. Em todos os casos, a
independência em relação a $R$ tem a mesma origem: potência e duração variam
inversamente.}
\end{questao}

\begin{questao}[4]{}
\textbf{(Projeto de banco.)} Deseja-se um banco de $\SI{100}{\micro\farad}$
capaz de operar em $\SI{100}{\volt}$, dispondo apenas de capacitores de
$\SI{100}{\micro\farad}$ e $\SI{50}{\volt}$.
\begin{itensq}
\item Determine a associação e o número de unidades.
\item Explique por que são necessários resistores de equalização.
\item Dimensione-os para uma corrente de fuga de até
      $\SI{10}{\micro\ampere}$.
\end{itensq}
\resposta{(a) duas séries de dois, em paralelo --- 4 unidades;
(c) $R\le\SI{50}{\kilo\ohm}$, dissipando $\SI{50}{\milli\watt}$ cada}
\resolucao{(a) \textbf{Série de dois:} suporta $\SI{100}{\volt}$, mas a
capacitância cai para $\SI{50}{\micro\farad}$.\\
\textbf{Duas dessas séries em paralelo:} $50+50=\SI{100}{\micro\farad}$,
mantendo $\SI{100}{\volt}$. Total: \textbf{4 capacitores}.\\
Note o custo: quadruplicou-se o número de componentes para dobrar a tensão. É o
preço inevitável --- a energia armazenada, essa sim, quadruplicou.\\
(b) \textbf{Sem equalização, a divisão de tensão em série é instável.} A tensão
de repouso é fixada pelas \emph{correntes de fuga}, não pelas capacitâncias: o
capacitor de menor fuga assume mais tensão. Como a fuga cresce com a tensão e
com a temperatura, o desequilíbrio se realimenta e pode levar um dos capacitores
a ultrapassar $\SI{50}{\volt}$ e perfurar --- o que joga toda a tensão no outro
e o destrói em seguida.\\
(c) Ligue um resistor em paralelo com cada capacitor, dimensionado para que sua
corrente domine a fuga. Adotando um fator $100$:
\[
I_R=100\cdot\pot{10}{-6}=\SI{1}{\milli\ampere}
\;\Longrightarrow\;
R\le\frac{50}{\pot{1}{-3}}=\SI{50}{\kilo\ohm}.
\]
Dissipação: $P=\dfrac{50^{2}}{\pot{50}{3}}=\SI{50}{\milli\watt}$ por resistor
--- especifique $\frac14\,\si{\watt}$.\\
\textbf{Bônus:} esses resistores também funcionam como \emph{bleeders},
descarregando o banco quando o equipamento é desligado. Com
$\tau=RC=50\,\mathrm{k}\times100\,\mu=\SI{5}{\second}$, o banco fica seguro em
menos de meio minuto.}
\end{questao}

\begin{questao}[4]{}
\textbf{(Projeto de filtro.)} Projete um passa-baixas RC de primeira ordem com
$f_c=\SI{100}{\hertz}$, usando $C=\SI{100}{\nano\farad}$.
\begin{itensq}
\item Determine $R$ e escolha o valor comercial E24.
\item Recalcule $f_c$ e avalie o erro.
\item Determine a atenuação em $\SI{1}{\kilo\hertz}$ e discuta a limitação da
      primeira ordem.
\end{itensq}
\resposta{(a) $R=\SI{15.9}{\kilo\ohm}\to\SI{16}{\kilo\ohm}$;
(b) $f_c=\SI{99.5}{\hertz}$, erro $0{,}5\%$;
(c) $\approx-\SI{20}{\decibel}$}
\resolucao{(a)
\[
R=\frac{1}{2\pi f_cC}=\frac{1}{2\pi(100)(\pot{100}{-9})}
=\SI{15915}{\ohm}.
\]
O valor E24 mais próximo é $\SI{16}{\kilo\ohm}$.\\
(b) $f_c=\dfrac{1}{2\pi(16000)(\pot{100}{-9})}=\SI{99.5}{\hertz}$ --- erro de
$0{,}5\%$, muito abaixo da tolerância do próprio capacitor (tipicamente
$\pm10\%$ ou pior). \textbf{O capacitor, e não o resistor, domina a incerteza
do projeto.}\\
(c) Uma década acima de $f_c$, o ganho é
\[
|H|=\frac{1}{\sqrt{1+(f/f_c)^{2}}}\approx\frac{1}{10}
\;\Longrightarrow\;
-\SI{20}{\decibel}.
\]
\textbf{Limitação:} $\SI{20}{\decibel}$ por década é pouco. Para atenuar
$\SI{60}{\decibel}$ seria preciso $f=\SI{100}{\kilo\hertz}$ --- três décadas
acima. Se a especificação exigir corte mais abrupto, é necessário aumentar a
\emph{ordem} do filtro (cascatear seções ou usar topologia ativa), não
apertar $R$ e $C$.\\
\textbf{Cuidado ao cascatear:} duas seções RC ligadas diretamente
\emph{carregam} uma à outra e o $f_c$ resultante não é o previsto. É preciso
desacoplá-las com um buffer, ou projetar a rede como um todo.}
\end{questao}

\begin{questao}[4]{}
\textbf{(Propagação de incerteza.)} Analise o efeito das tolerâncias de $R$ e
$C$ sobre a constante de tempo.
\begin{itensq}
\item Obtenha a expressão da incerteza relativa de $\tau$.
\item Avalie para $R$ com $\pm1\%$ e $C$ eletrolítico com $\pm20\%$, nos
      critérios de pior caso e estatístico.
\item Conclua sobre onde investir em precisão.
\end{itensq}
\resposta{(a) $\dfrac{\Delta\tau}{\tau}=\dfrac{\Delta R}{R}+\dfrac{\Delta C}{C}$;
(b) $21\%$ (pior caso) e $20{,}0\%$ (estatístico)}
\resolucao{(a) Como $\tau=RC$ é um \emph{produto}, tome o logaritmo e
diferencie:
\[
\ln\tau=\ln R+\ln C
\;\Longrightarrow\;
\frac{\Delta\tau}{\tau}=\frac{\Delta R}{R}+\frac{\Delta C}{C}.
\]
As incertezas \textbf{relativas} se somam --- diferentemente da soma, em que se
somam as absolutas.\\
(b) \emph{Pior caso:} $1\%+20\%=21\%$.\\
\emph{Estatístico:} $\sqrt{1^{2}+20^{2}}=\sqrt{401}=20{,}02\%$.\\
(c) \textbf{Os dois critérios praticamente coincidem} --- e isso é a
informação relevante. Quando uma parcela domina de forma tão desproporcional,
combinar em quadratura não ajuda: o resultado é essencialmente a maior
tolerância isolada.\\
\textbf{Conclusão de projeto:} trocar o resistor de $1\%$ por um de $0{,}1\%$
não muda nada ($20{,}00\%$ contra $20{,}02\%$). Todo o ganho está no capacitor.
Se a aplicação exigir temporização precisa, é preciso abandonar o eletrolítico
--- filme de poliéster oferece $\pm5\%$, e polipropileno chega a $\pm1\%$ com
excelente estabilidade térmica.\\
\textbf{Regra geral:} identifique a maior contribuição \emph{antes} de
especificar componentes. Precisão gasta em parcelas menores é dinheiro
jogado fora.}
\end{questao}

\begin{questao}[4]{}
\textbf{(Dedução --- força entre as placas.)} Considere um capacitor plano de
área $A$ e separação $d$.
\begin{itensq}
\item Obtenha a força entre as placas pelo método da energia, com carga
      constante.
\item Expresse o resultado em função de $U$ e discuta o caso a tensão
      constante.
\item Avalie para $A=\SI{100}{\centi\meter\squared}$,
      $d=\SI{1}{\milli\meter}$ e $U=\SI{1}{\kilo\volt}$.
\end{itensq}
\dados{$\varepsilon_0=\pot{8.85}{-12}\,\si{\farad\per\meter}$}
\resposta{$F=\dfrac{\varepsilon_0AU^{2}}{2d^{2}}\approx\SI{44}{\milli\newton}$,
atrativa}
\resolucao{(a) Com o capacitor \textbf{isolado}, $Q$ é constante e convém usar
\[
W=\frac{Q^{2}}{2C}=\frac{Q^{2}d}{2\varepsilon_0A}.
\]
A força que o sistema exerce é
\[
F=-\frac{\mathrm{d}W}{\mathrm{d}d}=-\frac{Q^{2}}{2\varepsilon_0A},
\]
negativa, isto é, \textbf{atrativa} --- as placas tendem a se aproximar,
reduzindo a energia.\\
(b) Substituindo $Q=CU=\dfrac{\varepsilon_0AU}{d}$:
\[
F=\frac{\varepsilon_0AU^{2}}{2d^{2}}\quad\text{(em módulo)}.
\]
\textbf{Sutileza importante:} a tensão constante, a energia \emph{aumenta} ao
aproximar as placas ($W=\frac12CU^{2}$ com $C$ crescendo), o que sugeriria
força repulsiva. O paradoxo se resolve incluindo a bateria: ela fornece
$U\,\mathrm{d}Q$, o dobro da variação de energia do campo. Feito o balanço
completo, a força é atrativa e de \emph{mesmo módulo} --- como tem de ser, pois
a força é uma grandeza física e não depende do que se escolhe manter
constante.\\
(c)
\[
F=\frac{(\pot{8.85}{-12})(\pot{1}{-2})(10^{6})}{2(10^{-6})}
=\SI{44.3}{\milli\newton}.
\]
\textbf{Ordem de grandeza:} cerca de $\SI{4.5}{\gram}$-força para
$\SI{1}{\kilo\volt}$ --- pequena, mas suficiente para fazer as placas vibrarem
audivelmente sob tensão alternada. É o "zumbido" de capacitores cerâmicos
classe 2, e o princípio de funcionamento de microfones e atuadores
eletrostáticos de MEMS.}
\end{questao}

\begin{questao}[4]{}
\textbf{(Projeto de temporizador.)} Um circuito deve atingir $70\%$ da tensão
final em exatamente $\SI{1}{\second}$.
\begin{itensq}
\item Determine $\tau$.
\item Escolha $C$ e determine $R$, com valores comerciais.
\item Avalie a repetibilidade real do temporizador.
\end{itensq}
\resposta{(a) $\tau=\SI{0.83}{\second}$; (b) $C=\SI{100}{\micro\farad}$ e
$R=\SI{8.2}{\kilo\ohm}$; (c) erro dominado pelo capacitor}
\resolucao{(a) De $0{,}70=1-e^{-t/\tau}$:
\[
e^{-t/\tau}=0{,}30
\;\Longrightarrow\;
\tau=\frac{t}{\ln(1/0{,}30)}=\frac{1}{1{,}204}=\SI{0.83}{\second}.
\]
(b) Com $C=\SI{100}{\micro\farad}$:
$R=\tau/C=0{,}83/\pot{100}{-6}=\SI{8306}{\ohm}$. O valor E24 mais próximo é
$\SI{8.2}{\kilo\ohm}$, que dá $\tau=\SI{0.82}{\second}$ e
\[
t_{70\%}=0{,}82(1{,}204)=\SI{0.987}{\second}\quad(\text{erro de }1{,}3\%).
\]
(c) \textbf{O erro de $1{,}3\%$ do arredondamento é irrelevante} diante da
tolerância do eletrolítico ($\pm20\%$), que se transfere integralmente para
$\tau$: o tempo real ficará entre $0{,}79$ e $\SI{1.19}{\second}$. Some-se a
deriva térmica (vários por cento entre $0$ e $\SI{60}{\celsius}$) e o
envelhecimento.\\
\textbf{Conclusão:} um temporizador RC com eletrolítico serve para "cerca de um
segundo", não para "um segundo". Para precisão, as alternativas são capacitor
de filme (levando a $R$ muito maior e $C$ menor) ou --- o caminho moderno ---
um contador digital com oscilador a cristal, cuja estabilidade é de partes por
milhão.\\
\textbf{Observação de método:} note que o requisito foi dado em $70\%$, e não
em $\tau$. Confundir os dois levaria a um erro de $20\%$ no projeto, pois
$t_{70\%}=1{,}204\,\tau$.}
\end{questao}

\begin{questao}[4]{}
\textbf{(Modelagem --- autodescarga.)} Um capacitor real é modelado por $C$
ideal em paralelo com uma resistência de fuga $R_f$.
\begin{itensq}
\item Deduza a lei de decaimento da tensão em circuito aberto.
\item Defina a constante de tempo de autodescarga e relacione-a com a
      especificação de "corrente de fuga" de catálogo.
\item Avalie para $C=\SI{1000}{\micro\farad}$ e
      $R_f=\SI{100}{\kilo\ohm}$, com tensão inicial de $\SI{300}{\volt}$.
\end{itensq}
\resposta{$u=U_0e^{-t/R_fC}$; $\tau=\SI{100}{\second}$; após
$\SI{5}{\minute}$ restam $\approx\SI{15}{\volt}$}
\resolucao{(a) Em circuito aberto, a única corrente é a de fuga:
\[
C\frac{\mathrm{d}u}{\mathrm{d}t}=-\frac{u}{R_f}
\;\Longrightarrow\;
u(t)=U_0e^{-t/R_fC}.
\]
(b) $\tau_f=R_fC$. Catálogos raramente informam $R_f$; informam a
\emph{corrente de fuga} $I_f$ a uma dada tensão. A conversão é imediata:
$R_f=U/I_f$, e portanto
\[
\tau_f=\frac{U\,C}{I_f}.
\]
Note que $\tau_f$ depende do \emph{produto} $R_fC$, e em eletrolíticos $R_f$
cai quando $C$ cresce --- razão pela qual é a autodescarga de capacitores grandes
não é proporcionalmente mais lenta.\\
(c) $\tau_f=\pot{100}{3}\cdot\pot{1000}{-6}=\SI{100}{\second}$. Após
$\SI{300}{\second}$:
\[
u=300\,e^{-3}=300(0{,}0498)=\SI{14.9}{\volt}.
\]
\textbf{Leitura de segurança:} cinco minutos depois de desligado, o barramento
ainda guarda $\SI{15}{\volt}$ --- e, o que é pior, $\SI{0.11}{\joule}$
concentrados. Aos $\SI{60}{\second}$ ainda havia $\SI{110}{\volt}$, plenamente
letais.\\
\textbf{Por isso o resistor de descarga é obrigatório}, e não opcional: com um
bleeder de $\SI{10}{\kilo\ohm}$, $\tau$ cai para $\SI{10}{\second}$ e o
equipamento fica seguro em menos de um minuto, ao custo de
$\SI{9}{\watt}$ apenas enquanto energizado.}
\end{questao}

\begin{questao}[4]{}
\textbf{(ESR e ondulação.)} Um capacitor de filtro de fonte chaveada tem
$\mathrm{ESR}=\SI{50}{\milli\ohm}$ e deve conduzir ondulação de
$\SI{2}{\ampere}$ eficazes.
\begin{itensq}
\item Determine a dissipação e a ondulação de tensão devida \emph{apenas} à
      ESR.
\item Compare com a ondulação devida à capacitância, para
      $C=\SI{470}{\micro\farad}$ a $\SI{100}{\kilo\hertz}$.
\item Conclua sobre o critério de seleção.
\end{itensq}
\resposta{(a) $\SI{0.2}{\watt}$ e $\approx\SI{100}{\milli\volt}$ de pico;
(b) $\approx\SI{5}{\milli\volt}$ --- a ESR domina}
\resolucao{(a) $P=\mathrm{ESR}\cdot I_{rms}^{2}=0{,}05(4)=\SI{0.2}{\watt}$.\\
A ondulação de tensão que a ESR produz acompanha a corrente instantaneamente:
\[
\Delta u_{ESR}\approx\mathrm{ESR}\cdot I_{pico}\approx0{,}05(2)=\SI{100}{\milli\volt}.
\]
(b) A parcela capacitiva depende da carga transferida por semiciclo. Tomando
$\Delta t\approx1/(2f)=\SI{5}{\micro\second}$:
\[
\Delta u_C\approx\frac{I\,\Delta t}{C}
=\frac{2(\pot{5}{-6})}{\pot{470}{-6}}\approx\SI{21}{\milli\volt}.
\]
(c) \textbf{A ESR domina} --- ela responde por cerca de $80\%$ da ondulação
total. Aumentar $C$ para $\SI{1000}{\micro\farad}$ reduziria $\Delta u_C$ pela
metade e deixaria a ondulação praticamente inalterada.\\
\textbf{Critério correto de seleção:} em fontes chaveadas de alta frequência,
escolhe-se o capacitor pela \textbf{ESR e pela corrente de ondulação
admissível}, não pela capacitância. É comum usar vários capacitores menores em
paralelo --- não para somar $C$, mas para dividir a ESR por $n$ e a dissipação
individual por $n^{2}$.\\
\textbf{Nota:} em frequências baixas ($\SI{120}{\hertz}$ de retificação de
rede), a relação se inverte e a capacitância volta a ser o critério dominante.
O parâmetro decisivo depende da faixa de operação.}
\end{questao}

\begin{questao}[4]{}
\textbf{(Comparação tecnológica.)} Um supercapacitor de $\SI{3000}{\farad}$ e
$\SI{2.7}{\volt}$ tem massa de $\SI{0.50}{\kilogram}$.
\begin{itensq}
\item Determine a energia armazenada e a densidade de energia.
\item Compare com uma bateria de íon-lítio ($\approx\SI{250}{\watt\hour\per\kilogram}$).
\item Aponte em que aspecto o supercapacitor é superior e por quê.
\end{itensq}
\resposta{(a) $\SI{10.9}{\kilo\joule}=\SI{3.04}{\watt\hour}$, ou
$\SI{6.1}{\watt\hour\per\kilogram}$; (b) cerca de $40$ vezes menos}
\resolucao{(a)
\[
W=\frac12CU^{2}=\frac12(3000)(2{,}7)^{2}=\SI{10935}{\joule}
=\frac{10935}{3600}=\SI{3.04}{\watt\hour}.
\]
Densidade: $3{,}04/0{,}50=\SI{6.1}{\watt\hour\per\kilogram}$.\\
(b) A bateria armazena cerca de \textbf{40 vezes mais} energia por quilograma.
Para igualar uma célula de íon-lítio de $\SI{20}{\gram}$, seriam necessários
$\SI{0.8}{\kilogram}$ de supercapacitores.\\
(c) \textbf{Potência específica e vida útil.} O supercapacitor pode ser
carregado e descarregado em \emph{segundos}, com correntes de centenas de
ampères, porque armazena energia \textbf{eletrostaticamente} --- sem reações
químicas, sem difusão iônica lenta, sem degradação de eletrodos. Suporta
centenas de milhares de ciclos, contra alguns milhares da bateria, e funciona
bem a baixas temperaturas.\\
\textbf{Consequência de aplicação:} baterias são \emph{reservatórios} de
energia; supercapacitores são \emph{amortecedores} de potência. Por isso
convivem em veículos elétricos --- o supercapacitor absorve o pico da frenagem
regenerativa e entrega o pico da aceleração, poupando a bateria justamente dos
regimes que mais a degradam.\\
\textbf{Ressalva:} a tensão do supercapacitor cai linearmente com a carga
(pois $Q=CU$), ao contrário do platô da bateria. Aproveitar mais de $75\%$ da
energia exige um conversor CC-CC de entrada ampla.}
\end{questao}

% END BANCO COMPLEMENTAR

\gabarito[10.1 Capacitores]

% =====================================================================
\subsection{Indutores}
% =====================================================================

\begin{conceito}[Antes de começar]
\textbf{Indutância de um solenoide:}
$L = \dfrac{\mu_0\,N^2\,A}{\ell}$, com
$\mu_0 = 4\pi\times10^{-7}\,\si{\henry\per\meter}$.\par
\textbf{Tensão induzida:} $v_L = L\,\dfrac{\mathrm{d}i}{\mathrm{d}t}$ --- o
indutor se opõe à \emph{variação} de corrente.\par
\textbf{Energia armazenada:} $E = \tfrac12 L\,I^2$.\par
\textbf{Associação (igual à dos resistores):} série soma, paralelo pelo inverso.\par
\textbf{Circuito RL:} $\tau = L/R$; na energização,
$i(t) = \dfrac{U_0}{R}\big(1 - e^{-t/\tau}\big)$.
\end{conceito}

\legendaniveis

\blocofixacao

\begin{questao}[1]{}
Um indutor ideal, em regime permanente CC (corrente constante), comporta-se como:
\begin{alternativas}
\item um circuito aberto.
\item um curto-circuito.
\item um resistor de valor $L$.
\item uma fonte de tensão.
\item um capacitor.
\end{alternativas}
\resposta{(b)}
\resolucao{Como $v_L = L\dfrac{\mathrm{d}i}{\mathrm{d}t}$ e, em regime permanente,
a corrente é constante ($\mathrm{d}i/\mathrm{d}t = 0$), a tensão no indutor é
\emph{nula}: ele se comporta como um \textbf{curto-circuito}. (É o oposto do
capacitor, que vira circuito aberto.)}
\end{questao}

\begin{questao}[1]{}
Calcule a indutância de um solenoide com $N = 500$ espiras, comprimento
$\ell = \SI{0.3}{\meter}$ e área $A = \SI{5}{\centi\meter\squared}$.
\dados{$\mu_0 = 4\pi\times10^{-7}\,\si{\henry\per\meter}$.}
\resposta{$L \approx \SI{0.52}{\milli\henry}$}
\resolucao{Com $A = \SI{5}{\centi\meter\squared} = \pot{5}{-4}\,\si{\meter\squared}$:
\[
L=\frac{\mu_0 N^2 A}{\ell}
=\frac{(\pot{4\pi}{-7})(500)^2(\pot{5}{-4})}{0{,}3}
\approx\pot{5.24}{-4}\,\si{\henry}=\SI{0.52}{\milli\henry}.
\]}
\end{questao}

\begin{questao}[1]{}
Um indutor de $\SI{0.1}{\henry}$ tem a corrente variando de $0$ a
$\SI{2}{\ampere}$ em $\SI{0.2}{\second}$. Calcule a tensão média induzida.
\resposta{$v_L = \SI{1}{\volt}$}
\resolucao{$v_L = L\dfrac{\Delta i}{\Delta t}
= 0{,}1\cdot\dfrac{2}{0{,}2} = \SI{1}{\volt}$.}
\end{questao}

\begin{questao}[1]{}
Dois indutores de $\SI{3}{\henry}$ e $\SI{9}{\henry}$ são ligados em paralelo. A
indutância equivalente é:
\begin{alternativas}[3]
\item \SI{12}{\henry}
\item \SI{6}{\henry}
\item \SI{2.25}{\henry}
\item \SI{3}{\henry}
\item \SI{27}{\henry}
\end{alternativas}
\resposta{(c)}
\resolucao{Indutores seguem a mesma regra dos resistores. Em paralelo:
\[
\Leq=\frac{L_1 L_2}{L_1+L_2}=\frac{3\cdot9}{12}=\frac{27}{12}=\SI{2.25}{\henry}.
\]}
\end{questao}

\blocoaplicacao

\begin{questao}[2]{}
Um indutor de $L = \SI{1}{\henry}$ é ligado em série com $R = \SI{100}{\ohm}$ a
uma fonte de $\SI{12}{\volt}$.
\begin{itensq}
\item Qual a constante de tempo?
\item Qual a corrente final (regime permanente)?
\item Qual a corrente após um intervalo igual a $\tau$?
\end{itensq}
\resposta{(a) $\tau = \SI{10}{\milli\second}$; (b) \SI{0.12}{\ampere};
(c) $\approx\SI{75.8}{\milli\ampere}$}
\resolucao{(a) $\tau = \dfrac{L}{R} = \dfrac{1}{100} = \SI{10}{\milli\second}$.\\
(b) Em regime permanente, o indutor vira curto e a corrente é limitada só por $R$:
$I_\infty = \dfrac{12}{100} = \SI{0.12}{\ampere}$.\\
(c) $i(\tau) = I_\infty(1 - e^{-1}) = 0{,}12\cdot0{,}632
\approx\SI{75.8}{\milli\ampere}$ --- cerca de $\SI{63}{\percent}$ do valor final,
o mesmo comportamento do capacitor na carga.}
\end{questao}

\begin{questao}[2]{}
Calcule a indutância equivalente para:
\begin{itensq}
\item $L_1 = \SI{2}{\milli\henry}$ e $L_2 = \SI{3}{\milli\henry}$ em série;
\item os mesmos dois indutores em paralelo.
\end{itensq}
\resposta{(a) \SI{5}{\milli\henry}; (b) \SI{1.2}{\milli\henry}}
\resolucao{(a) Série: $\Leq = L_1 + L_2 = 2 + 3 = \SI{5}{\milli\henry}$.\\
(b) Paralelo: $\Leq = \dfrac{2\cdot3}{5} = \SI{1.2}{\milli\henry}$.}
\end{questao}

\begin{questao}[2]{Apostila original revisada}
Dois indutores ideais de $\SI{5}{\henry}$ são ligados em paralelo. Essa associação
é conectada em série com um indutor de $\SI{2}{\henry}$. Determine a indutância
equivalente.
\resposta{$L_{\mathrm{eq}}=\SI{4.5}{\henry}$}
\resolucao{Os dois indutores iguais em paralelo resultam em
$L_p=5/2=\SI{2.5}{\henry}$. Em série com $\SI{2}{\henry}$:
$L_{\mathrm{eq}}=2{,}5+2=\SI{4.5}{\henry}$.}
\end{questao}

\blocoaprofundamento

\begin{questao}[3]{}
Compare, lado a lado, o comportamento de um capacitor e de um indutor em um
circuito de corrente contínua, preenchendo mentalmente a tabela abaixo e
justificando cada linha.

\begin{table}[H]
\centering
\small\sffamily
\setlength{\arraystretch}{1.3}
\begin{tabular}{@{}l c c@{}}
\toprule
 & \textbf{Capacitor} & \textbf{Indutor}\\
\midrule
Grandeza que não muda bruscamente & tensão $v_C$ & corrente $i_L$\\
Em regime permanente CC & circuito aberto & curto-circuito\\
Energia armazenada & $\tfrac12 C v^2$ (campo elétrico) & $\tfrac12 L i^2$ (campo magnético)\\
Constante de tempo & $\tau = RC$ & $\tau = L/R$\\
\bottomrule
\end{tabular}
\caption{Dualidade entre capacitor e indutor.}
\end{table}

Explique por que a tensão no capacitor e a corrente no indutor não podem variar
instantaneamente.
\resposta{Ver comentário}
\resolucao{A tensão do capacitor está ligada à sua carga por $v = Q/C$; mudar $v$
instantaneamente exigiria corrente infinita ($i = C\,\mathrm{d}v/\mathrm{d}t$),
o que é fisicamente impossível. De modo \emph{dual}, a corrente do indutor não
pode saltar, pois isso exigiria tensão infinita
($v = L\,\mathrm{d}i/\mathrm{d}t$). Por isso o capacitor "segura" a tensão e o
indutor "segura" a corrente --- é essa continuidade que dá origem às respostas
exponenciais com constante de tempo $\tau$. A tabela evidencia a bela
\textbf{dualidade} entre os dois componentes: trocando tensão$\leftrightarrow$corrente
e $C\leftrightarrow L$, uma descrição vira a outra.}
\end{questao}

\begin{questao}[3]{}
Um indutor de $L = \SI{0.5}{\henry}$ conduz corrente estabilizada de
$\SI{4}{\ampere}$, alimentado por uma fonte através de um resistor. Em
$t=0$, a chave que liga a fonte é \textbf{aberta bruscamente}.
\begin{itensq}
\item Qual a energia armazenada no indutor imediatamente antes da abertura?
\item Por que surge um arco elétrico (faísca) nos contatos da chave?
\item Estime a tensão induzida se a corrente for interrompida em
      $\SI{1}{\milli\second}$.
\item Qual solução prática se usa para proteger o circuito? Explique seu
      funcionamento.
\end{itensq}
\resposta{(a) $E=\SI{4}{\joule}$; (c) $\approx\SI{2000}{\volt}$; (b) e (d) ver comentário}
\resolucao{\textbf{(a)} $E = \tfrac12 L I^2 = \tfrac12(0{,}5)(4)^2
= \SI{4}{\joule}$ --- armazenados no \emph{campo magnético}.

\textbf{(b) Origem do arco.} A corrente em um indutor \textbf{não pode variar
instantaneamente}, pois isso exigiria tensão infinita:
$v_L = L\,\dfrac{\mathrm{d}i}{\mathrm{d}t}$. Ao abrir a chave, tenta-se impor
$\mathrm{d}i/\mathrm{d}t \to \infty$; o indutor reage gerando uma tensão
altíssima, de polaridade \emph{invertida}, que tenta manter a corrente
circulando. Essa tensão ioniza o ar entre os contatos e a corrente continua
fluindo pelo \textbf{arco} --- que é o caminho que o circuito "encontra" para
dissipar os \SI{4}{\joule} armazenados.

\textbf{(c) Estimativa da sobretensão.} Com $\Delta i = \SI{4}{\ampere}$ em
$\Delta t = \SI{1}{\milli\second}$:
\[
|v_L| = L\,\frac{\Delta i}{\Delta t}
= 0{,}5\cdot\frac{4}{\pot{1}{-3}} = \SI{2000}{\volt}.
\]
Duzentas vezes maior que uma alimentação típica de \SI{12}{\volt}! Se a abertura
fosse dez vezes mais rápida, a tensão seria dez vezes maior.

\textbf{(d) Proteção --- o diodo de roda livre (\emph{freewheeling}).}
Liga-se um diodo em \textbf{antiparalelo} com o indutor (catodo no positivo).
Em operação normal ele fica reversamente polarizado e não conduz. No instante da
abertura, a tensão do indutor inverte, o diodo entra em condução e oferece um
caminho fechado para a corrente, que decai suavemente com constante de tempo
$\tau = L/R$, dissipando a energia no próprio circuito em vez de no arco.\\
Alternativas: \emph{snubber} RC, varistor (MOV) ou diodo Zener --- este último
quando se deseja desmagnetização \emph{rápida}, aceitando uma sobretensão
controlada.\\
\textbf{Onde isso aparece na prática:} todo relé, contator, solenoide ou motor
CC acionado por transistor precisa desse diodo. Sem ele, o transistor de
acionamento é destruído na primeira comutação --- um dos erros mais comuns de
quem começa em eletrônica de potência.}
\end{questao}

\begin{questao}[3]{Apostila original revisada}
Um indutor de $L=\SI{2}{\milli\henry}$ permanece por muito tempo em série com um
resistor de $\SI{6}{\kilo\ohm}$ e uma fonte de $\SI{12}{\volt}$. Em $t=0$, a fonte
e o resistor de $\SI{6}{\kilo\ohm}$ são desconectados, e o indutor passa a descarregar
por um resistor de $\SI{2}{\kilo\ohm}$.
\begin{itensq}
\item Determine $i_L(0^-)$ e $i_L(0^+)$.
\item Calcule a constante de tempo da descarga.
\item Escreva $i_L(t)$ para $t\ge 0$ e determine o módulo da tensão inicial no resistor.
\end{itensq}
\resposta{$i_L(0^-)=i_L(0^+)=\SI{2}{\milli\ampere}$; $\tau=\SI{1}{\micro\second}$;
$i_L(t)=\SI{2}{\milli\ampere}e^{-t/(\SI{1}{\micro\second})}$; $|v_R(0^+)|=\SI{4}{\volt}$}
\resolucao{Em regime permanente, o indutor é curto-circuito: $i=12/6000=\SI{2}{\milli\ampere}$.
A corrente do indutor é contínua na comutação. Na descarga, $\tau=L/R=2\,\mathrm{mH}/2\,\mathrm{k\Omega}=\SI{1}{\micro\second}$.
A tensão inicial no resistor vale $Ri=2000\times0{,}002=\SI{4}{\volt}$; sua polaridade é tal que mantém a corrente.}
\end{questao}

% BEGIN BANCO COMPLEMENTAR Indutores
% =====================================================================
%  Banco complementar --- Indutores  (REVISADO)
%  Substitui a versao gerada por template (12 clones em N2, 12 em N3,
%  10 questoes conceituais indevidamente marcadas como N4).
%  Sem sobreposicao com o cap10, que ja traz: indutor ideal em CC (MC),
%  indutancia de solenoide N=500, rampa de corrente basica, paralelo
%  3H/9H, RL com L=1H e R=100 ohm, Leq serie/paralelo e a tabela
%  comparativa capacitor x indutor.
%  Distribuicao mantida: 6 N1 + 12 N2 + 12 N3 + 10 N4 = 40
% =====================================================================

\subsubsection*{Banco complementar --- Indutores}

\begin{atencao}[Organização do banco]
Três relações sustentam tudo: $v=L\dfrac{\mathrm{d}i}{\mathrm{d}t}$,
$W=\frac12LI^{2}$ e $\lambda=LI$. Delas seguem os dois fatos que governam o
comportamento transitório --- a corrente \textbf{não pode saltar} (isso exigiria
tensão infinita) e o indutor \textbf{se opõe à variação}, não à corrente em si.
O N3 trata de sobretensão, acoplamento, saturação e extração de parâmetros; o
N4 exige dedução da EDO, integração de energia e projeto com restrições
físicas.
\end{atencao}

\blocofixacao

\begin{questao}[1]{}
A tensão sobre um indutor vale $\SI{50}{\volt}$ quando a corrente varia à taxa
de $\SI{250}{\ampere\per\second}$. Determine a indutância.
\resposta{$L=\SI{0.2}{\henry}$}
\resolucao{Isolando $L$ na relação fundamental:
\[
v=L\frac{\mathrm{d}i}{\mathrm{d}t}
\;\Longrightarrow\;
L=\frac{v}{\mathrm{d}i/\mathrm{d}t}=\frac{50}{250}=\SI{0.2}{\henry}.
\]
O henry é, portanto, volt-segundo por ampère: $\SI{1}{\henry}$ é a indutância
que produz $\SI{1}{\volt}$ para uma variação de $\SI{1}{\ampere\per\second}$.}
\end{questao}

\begin{questao}[1]{}
Um indutor de $\SI{0.4}{\henry}$ conduz $\SI{5}{\ampere}$ constantes. Determine
a energia armazenada no campo magnético.
\resposta{$W=\SI{5}{\joule}$}
\resolucao{
\[
W=\frac12LI^{2}=\frac12(0{,}4)(5)^{2}=\SI{5}{\joule}.
\]
A dependência é \emph{quadrática} na corrente: dobrar $I$ quadruplica a
energia. Note que, embora a tensão sobre o indutor ideal seja nula em regime CC
(corrente constante), a energia armazenada não é --- ela ficou guardada no
campo durante a fase de crescimento da corrente.}
\end{questao}

\begin{questao}[1]{}
Sobre a corrente em um indutor ideal, é correto afirmar que ela:
\begin{alternativas}[2]
\item pode variar instantaneamente
\item não pode variar instantaneamente
\item é sempre constante
\item é sempre nula em corrente contínua
\end{alternativas}
\resposta{(b)}
\resolucao{Um salto de corrente significaria
$\mathrm{d}i/\mathrm{d}t\to\infty$ e, por $v=L\,\mathrm{d}i/\mathrm{d}t$, uma
tensão infinita --- fisicamente impossível. Logo a corrente no indutor é uma
função \textbf{contínua} do tempo: $i(0^{+})=i(0^{-})$.\\
Cuidado com (c) e (d): a corrente pode variar (e varia, nos transitórios), e em
CC ela é constante e em geral \emph{não nula}. O que é nula em CC é a
\emph{tensão}.}
\end{questao}

\begin{questao}[1]{}
Um indutor de $\SI{200}{\milli\henry}$ está em série com $\SI{50}{\ohm}$.
Determine a constante de tempo do circuito.
\resposta{$\tau=\SI{4}{\milli\second}$}
\resolucao{
\[
\tau=\frac{L}{R}=\frac{0{,}200}{50}=\SI{4e-3}{\second}=\SI{4}{\milli\second}.
\]
\textbf{Atenção à estrutura:} no circuito RL, $\tau=L/R$ --- resistência
\emph{maior} torna o transitório \emph{mais rápido}. No RC é o oposto
($\tau=RC$). Trocar as duas fórmulas é o erro mais comum entre os dois
elementos armazenadores.}
\end{questao}

\begin{questao}[1]{}
Determine o fluxo concatenado em um indutor de $\SI{0.25}{\henry}$ percorrido
por $\SI{8}{\ampere}$.
\resposta{$\lambda=\SI{2}{\weber}$}
\resolucao{
\[
\lambda=LI=0{,}25\cdot8=\SI{2}{\weber}.
\]
O fluxo concatenado $\lambda=N\Phi$ é o produto do fluxo por espira pelo número
de espiras. A indutância é justamente a constante de proporcionalidade entre
$\lambda$ e $I$ --- é essa a definição primária de $L$, da qual
$v=L\,\mathrm{d}i/\mathrm{d}t$ decorre pela lei de Faraday.}
\end{questao}

\begin{questao}[1]{}
Aplicam-se $\SI{30}{\volt}$ a um indutor de $\SI{0.15}{\henry}$ inicialmente
sem corrente. Determine a taxa inicial de variação da corrente.
\resposta{$\mathrm{d}i/\mathrm{d}t=\SI{200}{\ampere\per\second}$}
\resolucao{
\[
\frac{\mathrm{d}i}{\mathrm{d}t}=\frac{v}{L}=\frac{30}{0{,}15}
=\SI{200}{\ampere\per\second}.
\]
No instante inicial a corrente é nula, mas sua \emph{taxa} é máxima. Se a
tensão se mantivesse constante e não houvesse resistência, a corrente cresceria
indefinidamente em rampa --- é a resistência do circuito que estabelece o
patamar final.}
\end{questao}

\blocoaplicacao

\begin{questao}[2]{}
Um indutor de $\SI{0.5}{\henry}$ conduz $\SI{3}{\ampere}$.
\begin{itensq}
\item Determine a energia armazenada.
\item Determine a tensão necessária para levar a corrente a zero
      \emph{linearmente} em $\SI{50}{\milli\second}$.
\end{itensq}
\resposta{(a) $\SI{2.25}{\joule}$; (b) $\SI{30}{\volt}$}
\resolucao{(a) $W=\frac12(0{,}5)(9)=\SI{2.25}{\joule}$.\\
(b) A taxa é $\mathrm{d}i/\mathrm{d}t=(0-3)/0{,}050
=\SI{-60}{\ampere\per\second}$, e
\[
|v|=L\left|\frac{\mathrm{d}i}{\mathrm{d}t}\right|=0{,}5\cdot60=\SI{30}{\volt}.
\]
O sinal negativo indica \textbf{inversão de polaridade}: o indutor passa a
funcionar como fonte, devolvendo ao circuito os $\SI{2.25}{\joule}$
armazenados. É esse comportamento que produz as sobretensões de abertura.}
\end{questao}

\begin{questao}[2]{}
A corrente em um indutor de $\SI{100}{\milli\henry}$ cai linearmente de
$\SI{4}{\ampere}$ a zero em $\SI{20}{\milli\second}$. Determine a tensão
induzida e sua polaridade.
\resposta{$\SI{20}{\volt}$, com polaridade \emph{invertida} em relação à fase
de crescimento}
\resolucao{
\[
|v|=L\left|\frac{\Delta i}{\Delta t}\right|
=0{,}100\cdot\frac{4}{0{,}020}=\SI{20}{\volt}.
\]
Durante o \emph{crescimento} da corrente, o indutor absorve energia e sua
tensão se opõe ao aumento. Durante o \emph{decrescimento}, ele devolve energia
e a tensão inverte, tentando sustentar a corrente.\\
Em ambos os casos vale a lei de Lenz: o indutor se opõe à \emph{variação}, não
ao sentido da corrente.}
\end{questao}

\begin{questao}[2]{}
Três indutores não acoplados de $\SI{6}{\milli\henry}$,
$\SI{12}{\milli\henry}$ e $\SI{12}{\milli\henry}$ são associados. Determine
$L_{eq}$ em série e em paralelo.
\resposta{Série: $\SI{30}{\milli\henry}$; paralelo: $\SI{3}{\milli\henry}$}
\resolucao{\textbf{Série:} $L_{eq}=6+12+12=\SI{30}{\milli\henry}$.\\
\textbf{Paralelo:}
\[
\frac{1}{L_{eq}}=\frac16+\frac1{12}+\frac1{12}=\frac{2+1+1}{12}=\frac13
\;\Longrightarrow\;
L_{eq}=\SI{3}{\milli\henry}.
\]
As regras são \emph{idênticas} às dos resistores --- e opostas às dos
capacitores. A ressalva é a hipótese de \textbf{não acoplamento}: se houver
fluxo mútuo, as fórmulas mudam (ver bloco de aprofundamento).}
\end{questao}

\begin{questao}[2]{}
Um circuito RL série tem $L=\SI{0.6}{\henry}$, $R=\SI{30}{\ohm}$ e fonte de
$\SI{24}{\volt}$, com a chave fechada em $t=0$.
\begin{itensq}
\item Determine $\tau$ e a corrente final.
\item Determine a tensão sobre o indutor em $t=0^{+}$.
\end{itensq}
\resposta{(a) $\tau=\SI{20}{\milli\second}$, $I_f=\SI{0.8}{\ampere}$;
(b) $\SI{24}{\volt}$}
\resolucao{(a) $\tau=L/R=0{,}6/30=\SI{20}{\milli\second}$;
$I_f=V/R=24/30=\SI{0.8}{\ampere}$.\\
(b) Em $t=0^{+}$ a corrente ainda é nula (continuidade), logo o resistor não
tem queda e \textbf{toda} a tensão da fonte aparece sobre o indutor:
$v_L(0^{+})=\SI{24}{\volt}$.\\
\textbf{Os dois extremos:} em $t=0^{+}$ o indutor se comporta como circuito
\emph{aberto} (corrente nula); em $t\to\infty$, como \emph{curto}
(tensão nula). Reconhecer esses dois limites resolve a maioria dos problemas
sem integrar nada.}
\end{questao}

\begin{questao}[2]{}
No circuito anterior, determine a corrente em $t=\tau$, $t=2\tau$ e
$t=3\tau$.
\resposta{$\SI{0.506}{\ampere}$, $\SI{0.692}{\ampere}$ e
$\SI{0.760}{\ampere}$}
\resolucao{Com $i(t)=I_f\left(1-e^{-t/\tau}\right)$ e
$I_f=\SI{0.8}{\ampere}$:
\[
i(\tau)=0{,}8(1-e^{-1})=0{,}8(0{,}632)=\SI{0.506}{\ampere},
\]
\[
i(2\tau)=0{,}8(0{,}865)=\SI{0.692}{\ampere},
\qquad
i(3\tau)=0{,}8(0{,}950)=\SI{0.760}{\ampere}.
\]
As frações $63{,}2\%$, $86{,}5\%$ e $95{,}0\%$ são universais --- valem para
qualquer RL ou RC. Após $5\tau$ atinge-se $99{,}3\%$, e por convenção o
transitório é dado por encerrado.}
\end{questao}

\begin{questao}[2]{}
Aplica-se a um indutor de $\SI{5}{\milli\henry}$ uma onda quadrada de tensão
que alterna entre $+\SI{10}{\volt}$ e $-\SI{10}{\volt}$, com meio período de
$\SI{1}{\milli\second}$. Determine a variação de corrente pico a pico em
regime.
\resposta{$\Delta i=\SI{2}{\ampere}$}
\resolucao{Durante o semiciclo positivo, a tensão constante produz rampa de
corrente:
\[
\Delta i=\frac{v\,\Delta t}{L}=\frac{10\cdot0{,}001}{0{,}005}=\SI{2}{\ampere}.
\]
No semiciclo negativo a corrente desce pela mesma quantidade, e o regime é
simétrico.\\
\textbf{Resultado geral:} tensão quadrada no indutor produz corrente
\textbf{triangular} --- o indutor \emph{integra} a tensão. O capacitor faz o
oposto: corrente quadrada nele produz tensão triangular.}
\end{questao}

\begin{questao}[2]{}
Um indutor \emph{real} de $\SI{0.5}{\henry}$ possui resistência de enrolamento
$r=\SI{8}{\ohm}$ e é ligado a uma fonte CC de $\SI{24}{\volt}$.
\begin{itensq}
\item Determine a corrente em regime permanente.
\item Determine a energia armazenada e a potência dissipada.
\end{itensq}
\resposta{(a) $\SI{3}{\ampere}$; (b) $W=\SI{2.25}{\joule}$ e
$P=\SI{72}{\watt}$}
\resolucao{(a) Em regime, $\mathrm{d}i/\mathrm{d}t=0$ e o indutor ideal não cai
tensão: só $r$ limita a corrente.
\[
I=\frac{24}{8}=\SI{3}{\ampere}.
\]
(b) $W=\frac12(0{,}5)(9)=\SI{2.25}{\joule}$;
$P=rI^{2}=8(9)=\SI{72}{\watt}$.\\
\textbf{Contraste essencial:} a energia é \emph{armazenada} e recuperável; a
potência é \emph{dissipada} continuamente e perdida. Manter esses
$\SI{2.25}{\joule}$ no campo custa $\SI{72}{\watt}$ permanentes --- razão pela
qual eletroímãs de grande porte usam supercondutores, em que $r=0$ e a
manutenção do campo é gratuita.}
\end{questao}

\begin{questao}[2]{}
Dois indutores armazenam a mesma energia: o primeiro tem
$L_1=\SI{2}{\henry}$ com $I_1=\SI{1}{\ampere}$. Se
$L_2=\SI{0.5}{\henry}$, determine $I_2$.
\resposta{$I_2=\SI{2}{\ampere}$}
\resolucao{
\[
\frac12L_1I_1^{2}=\frac12L_2I_2^{2}
\;\Longrightarrow\;
I_2=I_1\sqrt{\frac{L_1}{L_2}}=1\cdot\sqrt{4}=\SI{2}{\ampere}.
\]
Ambos armazenam $\SI{1}{\joule}$. Reduzir $L$ por um fator $4$ exige
\emph{dobrar} a corrente --- e, como as perdas vão com $I^{2}$, o indutor menor
custa quatro vezes mais em dissipação para armazenar o mesmo tanto. É o
compromisso central no projeto de indutores de potência.}
\end{questao}

\begin{questao}[2]{}
Um solenoide de $N=200$ espiras tem indutância $L$. Determine a nova
indutância se o número de espiras for dobrado, mantidos o comprimento e a
seção.
\resposta{$4L$}
\resolucao{Para um solenoide longo,
\[
L=\frac{\mu N^{2}A}{\ell},
\]
de modo que $L\propto N^{2}$. Dobrar $N$ quadruplica $L$.\\
\textbf{Por que o quadrado:} dobrar as espiras dobra o campo produzido por
uma dada corrente \emph{e} dobra o número de espiras que concatenam esse campo.
Os dois efeitos se multiplicam.\\
\textbf{Ressalva prática:} com o mesmo espaço de bobina, dobrar $N$ exige fio
de seção menor, o que aumenta $r$ e as perdas. Na prática, o ganho de $L$ vem
acompanhado de penalidade em corrente admissível.}
\end{questao}

\begin{questao}[2]{}
Um indutor com núcleo de ar tem $L_0=\SI{2}{\milli\henry}$. Introduz-se um
núcleo ferromagnético de permeabilidade relativa $\mu_r=500$. Determine a nova
indutância e comente as limitações.
\resposta{$L=\SI{1}{\henry}$}
\resolucao{Como $L\propto\mu=\mu_r\mu_0$:
\[
L=\mu_rL_0=500\cdot2=\SI{1000}{\milli\henry}=\SI{1}{\henry}.
\]
\textbf{Três limitações do núcleo.} \emph{(i)} \textbf{Saturação:} acima de
certa corrente, $\mu_r$ despenca e $L$ colapsa. \emph{(ii)} \textbf{Perdas:}
histerese e correntes de Foucault dissipam energia proporcionalmente à
frequência. \emph{(iii)} \textbf{Não linearidade:} $L$ deixa de ser constante,
e as fórmulas lineares passam a ser aproximações válidas apenas em torno de um
ponto de operação.\\
Núcleos de ar têm $L$ minúsculo, mas são perfeitamente lineares e sem perdas
magnéticas --- por isso dominam aplicações de radiofrequência.}
\end{questao}

\begin{questao}[2]{}
No circuito RL com $\tau=\SI{20}{\milli\second}$ e $I_f=\SI{0.8}{\ampere}$,
determine o instante em que a corrente atinge $\SI{0.5}{\ampere}$.
\resposta{$t\approx\SI{19.6}{\milli\second}$}
\resolucao{De $i=I_f(1-e^{-t/\tau})$, isole a exponencial:
\[
e^{-t/\tau}=1-\frac{0{,}5}{0{,}8}=0{,}375
\;\Longrightarrow\;
t=\tau\ln\!\left(\frac{1}{0{,}375}\right)=0{,}020\cdot0{,}981
=\SI{19.6}{\milli\second}.
\]
Ou seja, praticamente $1\tau$ --- coerente com o fato de que
$\SI{0.5}{\ampere}$ corresponde a $62{,}5\%$ de $I_f$, muito próximo dos
$63{,}2\%$ atingidos em exatamente $\tau$.}
\end{questao}

\begin{questao}[2]{}
Um indutor de $\SI{0.4}{\henry}$ conduz $\SI{2}{\ampere}$ quando a fonte é
substituída por um curto através de $R=\SI{20}{\ohm}$. Determine $\tau$ e a
corrente após $\SI{40}{\milli\second}$.
\resposta{$\tau=\SI{20}{\milli\second}$; $i=\SI{0.27}{\ampere}$}
\resolucao{$\tau=L/R=0{,}4/20=\SI{20}{\milli\second}$.\\
Na descarga, a corrente decai exponencialmente a partir do valor inicial:
\[
i(t)=I_0e^{-t/\tau}=2\,e^{-40/20}=2e^{-2}=2(0{,}135)=\SI{0.27}{\ampere}.
\]
\textbf{Continuidade:} no instante da comutação a corrente permanece
$\SI{2}{\ampere}$ --- ela não pode saltar. O que muda instantaneamente é a
\emph{tensão}, que assume $-RI_0=\SI{-40}{\volt}$ para manter essa corrente
circulando pelo resistor.}
\end{questao}

\blocoaprofundamento

\begin{questao}[3]{}
Um indutor de $\SI{0.2}{\henry}$ conduz $\SI{5}{\ampere}$ e é descarregado
através de um resistor.
\begin{itensq}
\item Determine a energia total dissipada no resistor.
\item Explique por que ela não depende do valor de $R$.
\item Determine como $R$ afeta o processo.
\end{itensq}
\resposta{(a) $\SI{2.5}{\joule}$; (b) conservação de energia; (c) $R$ altera a
\emph{duração} e a \emph{potência}, não a energia total}
\resolucao{(a) Toda a energia armazenada acaba no resistor:
\[
W=\frac12LI_0^{2}=\frac12(0{,}2)(25)=\SI{2.5}{\joule}.
\]
(b) O indutor ideal não dissipa e o circuito é fechado: a energia magnética não
tem outro destino. A independência em relação a $R$ é consequência direta da
conservação --- e é confirmada por integração no bloco de desafio.\\
(c) $R$ governa $\tau=L/R$ e, portanto, o \emph{ritmo}:
\begin{itemize}
\item $R$ pequeno $\Rightarrow$ $\tau$ grande: descarga lenta, potência baixa e
prolongada.
\item $R$ grande $\Rightarrow$ $\tau$ pequeno: descarga rápida, com pico de
potência $P_0=RI_0^{2}$ muito elevado.
\end{itemize}
\textbf{Exemplo:} com $R=\SI{10}{\ohm}$, $P_0=\SI{250}{\watt}$ por
$\SI{20}{\milli\second}$; com $R=\SI{1000}{\ohm}$, $P_0=\SI{25}{\kilo\watt}$
por $\SI{0.2}{\milli\second}$. Mesma energia, potências que diferem em cem
vezes.}
\end{questao}

\begin{questao}[3]{}
Um indutor de $\SI{0.1}{\henry}$ conduz $\SI{2}{\ampere}$ quando uma chave
mecânica é aberta, interrompendo a corrente em cerca de
$\SI{1}{\micro\second}$.
\begin{itensq}
\item Estime a tensão induzida.
\item Explique por que esse valor não se observa na prática.
\item Indique a consequência real.
\end{itensq}
\resposta{(a) $\SI{200}{\kilo\volt}$ (estimativa idealizada);
(b) o ar rompe muito antes; (c) formação de arco elétrico nos contatos}
\resolucao{(a)
\[
|v|=L\frac{\Delta i}{\Delta t}=0{,}1\cdot\frac{2}{10^{-6}}
=\pot{2}{5}\,\si{\volt}=\SI{200}{\kilo\volt}.
\]
(b) A rigidez dielétrica do ar é da ordem de
$\SI{3}{\kilo\volt\per\milli\meter}$. Com contatos separados por décimos de
milímetro, o ar ioniza a alguns \emph{quilovolts} --- muito antes dos
$\SI{200}{\kilo\volt}$. O cálculo é uma \textbf{estimativa por absurdo}: ele
mostra que a hipótese "a corrente cessa em $\SI{1}{\micro\second}$" é
insustentável.\\
(c) Forma-se um \textbf{arco elétrico}, que mantém a corrente circulando
enquanto o campo se esvazia. O arco limita a tensão a algumas centenas de
volts, mas corrói os contatos, gera interferência eletromagnética e reduz
drasticamente a vida da chave.\\
\textbf{Leitura correta:} não é o indutor que "gera" $\SI{200}{\kilo\volt}$ ---
é a impossibilidade de interromper a corrente que \emph{força} o circuito a
encontrar um caminho, e ele o encontra no ar.}
\end{questao}

\begin{questao}[3]{}
Um diodo de roda livre é colocado em antiparalelo com uma bobina de
$\SI{0.1}{\henry}$ e $r=\SI{5}{\ohm}$, que conduz $\SI{2}{\ampere}$.
\begin{itensq}
\item Explique o funcionamento no instante da abertura.
\item Determine a tensão máxima sobre a chave.
\item Determine a constante de tempo de extinção.
\end{itensq}
\resposta{(b) $\approx\SI{0.7}{\volt}$ acima da alimentação;
(c) $\tau=\SI{20}{\milli\second}$}
\resolucao{(a) Com a chave aberta, a corrente da bobina precisa de um caminho.
O diodo, que estava reversamente polarizado, entra em condução e fecha uma
malha bobina--diodo. A corrente continua circulando e decai exponencialmente,
dissipando-se em $r$ e no diodo --- sem arco.\\
(b) A tensão sobre a bobina fica grampeada em $-V_D\approx\SI{-0.7}{\volt}$, de
modo que a chave suporta apenas a tensão de alimentação acrescida de
$\SI{0.7}{\volt}$. Compare com os quilovolts do caso anterior.\\
(c) $\tau=L/r=0{,}1/5=\SI{20}{\milli\second}$.\\
\textbf{Custo do método:} a extinção fica \emph{lenta}. Em um relé, isso
prolonga o tempo de desarme; em acionamentos rápidos, prefere-se um diodo em
série com um resistor (ou um Zener), que eleva a tensão de grampeamento e
acelera a queda --- trocando proteção máxima por velocidade.}
\end{questao}

\begin{questao}[3]{}
Dois indutores acoplados, $L_1=\SI{4}{\milli\henry}$ e
$L_2=\SI{9}{\milli\henry}$, têm indutância mútua $M=\SI{3}{\milli\henry}$.
\begin{itensq}
\item Determine $L_{eq}$ em série nas duas orientações possíveis.
\item Determine o coeficiente de acoplamento $k$.
\item Explique como identificar experimentalmente a orientação.
\end{itensq}
\resposta{(a) $\SI{19}{\milli\henry}$ e $\SI{7}{\milli\henry}$; (b) $k=0{,}5$}
\resolucao{(a) Em série,
\[
L_{eq}=L_1+L_2\pm2M
\;\Longrightarrow\;
4+9+6=\SI{19}{\milli\henry}
\ \text{ou}\ 4+9-6=\SI{7}{\milli\henry}.
\]
O sinal $+$ vale para \textbf{acoplamento aditivo} (os fluxos se reforçam; a
corrente entra pelos dois pontos marcados); o sinal $-$, para
\textbf{subtrativo}.\\
(b)
\[
k=\frac{M}{\sqrt{L_1L_2}}=\frac{3}{\sqrt{36}}=0{,}5.
\]
Como $0\le k\le1$, este é um acoplamento médio --- típico de bobinas
próximas com núcleo de ar. Transformadores de potência atingem $k>0{,}99$.\\
(c) \textbf{Método experimental:} meça $L_{eq}$ nas duas ligações possíveis. A
\emph{maior} corresponde ao acoplamento aditivo e identifica os pontos.
Além disso, $M$ sai de graça:
\[
M=\frac{L_{aditivo}-L_{subtrativo}}{4}=\frac{19-7}{4}=\SI{3}{\milli\henry}.
\]
É assim que se determina $M$ sem acesso à geometria interna das bobinas.}
\end{questao}

\begin{questao}[3]{}
Uma fonte de $\SI{12}{\volt}$ alimenta, através de $R_1=\SI{20}{\ohm}$, um nó
do qual saem $R_3=\SI{60}{\ohm}$ ao terra e um ramo com
$R_2=\SI{30}{\ohm}$ em série com $L=\SI{0.9}{\henry}$.
\begin{itensq}
\item Determine a resistência de Thévenin vista pelo indutor.
\item Determine $\tau$ e a corrente final no indutor.
\end{itensq}
\resposta{(a) $R_{Th}=\SI{45}{\ohm}$; (b) $\tau=\SI{20}{\milli\second}$ e
$I_f=\SI{0.2}{\ampere}$}
\resolucao{(a) Zerando a fonte (curto), o indutor vê $R_2$ em série com
$R_1\parallel R_3$:
\[
R_{Th}=30+\frac{20\cdot60}{80}=30+15=\SI{45}{\ohm}.
\]
(b) $\tau=L/R_{Th}=0{,}9/45=\SI{20}{\milli\second}$.\\
A tensão de Thévenin sai com o indutor removido (sem corrente em $R_2$, logo
sem queda nele):
\[
V_{Th}=12\cdot\frac{60}{20+60}=\SI{9}{\volt}
\;\Longrightarrow\;
I_f=\frac{9}{45}=\SI{0.2}{\ampere}.
\]
\textbf{Método geral:} em qualquer circuito com \emph{um} elemento armazenador,
substitua todo o restante pelo equivalente de Thévenin visto por ele. O
problema se reduz a um RL série, com $\tau=L/R_{Th}$ e
$I_f=V_{Th}/R_{Th}$ --- e não é preciso resolver a rede inteira em regime
transitório.}
\end{questao}

\begin{questao}[3]{}
Um indutor de $\SI{0.5}{\henry}$ já conduz $\SI{1}{\ampere}$ quando é ligado a
uma fonte que estabeleceria $\SI{3}{\ampere}$ em regime, com
$R=\SI{25}{\ohm}$.
\begin{itensq}
\item Escreva $i(t)$.
\item Determine a corrente após $\SI{20}{\milli\second}$.
\item Determine a tensão inicial sobre o indutor.
\end{itensq}
\resposta{(a) $i(t)=3-2e^{-t/\tau}$ com $\tau=\SI{20}{\milli\second}$;
(b) $\SI{2.26}{\ampere}$; (c) $\SI{50}{\volt}$}
\resolucao{(a) A forma geral do transitório de primeira ordem é
\[
i(t)=I_f+\left(I_0-I_f\right)e^{-t/\tau},
\]
com $I_0=\SI{1}{\ampere}$, $I_f=\SI{3}{\ampere}$ e
$\tau=0{,}5/25=\SI{20}{\milli\second}$:
\[
i(t)=3-2e^{-t/0{,}020}.
\]
(b) $i(0{,}020)=3-2e^{-1}=3-0{,}736=\SI{2.26}{\ampere}$.\\
(c) Em $t=0^{+}$, $i=\SI{1}{\ampere}$ e o resistor cai
$25(1)=\SI{25}{\volt}$; a fonte vale $RI_f=25(3)=\SI{75}{\volt}$. Logo
\[
v_L(0^{+})=75-25=\SI{50}{\volt}.
\]
\textbf{Generalidade:} a fórmula de (a) cobre carga, descarga e condições
iniciais quaisquer --- basta identificar $I_0$, $I_f$ e $\tau$. Não é preciso
memorizar três expressões diferentes.}
\end{questao}

\begin{questao}[3]{}
Em um circuito RL, mediu-se corrente final de $\SI{2}{\ampere}$ com fonte de
$\SI{50}{\volt}$, e verificou-se que a corrente atingiu $\SI{1.264}{\ampere}$
em $\SI{8}{\milli\second}$. Determine $R$ e $L$.
\resposta{$R=\SI{25}{\ohm}$; $L=\SI{0.2}{\henry}$}
\resolucao{\textbf{Resistência:} do regime permanente,
$R=V/I_f=50/2=\SI{25}{\ohm}$.\\
\textbf{Constante de tempo:} note que
$1{,}264/2=0{,}632$ --- exatamente a fração característica de $1\tau$. Logo
$\tau=\SI{8}{\milli\second}$, e
\[
L=R\tau=25(0{,}008)=\SI{0.2}{\henry}.
\]
\textbf{Se a fração não fosse reconhecível}, o caminho geral seria
$\tau=-t/\ln(1-i/I_f)$.\\
\textbf{Método:} o regime permanente entrega $R$; o transitório entrega $\tau$
e, por consequência, $L$. Duas medidas de naturezas diferentes são necessárias,
porque nenhuma delas isola $L$ sozinha.}
\end{questao}

\begin{questao}[3]{}
Analise o sinal da potência instantânea $p=vi$ em um indutor durante um ciclo
completo de carga e descarga, e relacione-o com a energia armazenada.
\resposta{$p>0$ na carga (absorve), $p<0$ na descarga (devolve); a integral em
um ciclo completo é nula}
\resolucao{\textbf{Carga:} $i>0$ e crescente $\Rightarrow$
$v=L\,\mathrm{d}i/\mathrm{d}t>0$ $\Rightarrow$ $p>0$: o indutor \emph{absorve}
energia da fonte e a armazena no campo.\\
\textbf{Descarga:} $i>0$ mas decrescente $\Rightarrow$ $v<0$ $\Rightarrow$
$p<0$: o indutor \emph{fornece} energia ao circuito.\\
\textbf{Integral:} como
\[
\int p\,\mathrm{d}t=\int Li\,\frac{\mathrm{d}i}{\mathrm{d}t}\,\mathrm{d}t
=\int Li\,\mathrm{d}i=\frac12Li^{2},
\]
a energia depende \emph{apenas} da corrente instantânea, não do caminho
percorrido. Se o ciclo retorna à mesma corrente, a integral é nula: o indutor
ideal não consome energia líquida.\\
\textbf{Consequência:} indutor e capacitor são elementos \textbf{reativos} ---
trocam energia com o circuito sem consumi-la. Só a resistência dissipa. Essa
distinção é a base de todo o tratamento de potência ativa e reativa em corrente
alternada.}
\end{questao}

\begin{questao}[3]{}
Determine a densidade de energia magnética em um núcleo onde
$B=\SI{1.2}{\tesla}$, supondo $\mu\approx\mu_0$ no entreferro.
\begin{itensq}
\item Calcule $u$.
\item Compare com a densidade de energia elétrica em um campo de
      $\SI{3}{\mega\volt\per\meter}$.
\end{itensq}
\dados{$\mu_0=4\pi\times10^{-7}\,\si{\henry\per\meter}$;
$\varepsilon_0=\pot{8.85}{-12}\,\si{\farad\per\meter}$}
\resposta{(a) $u\approx\SI{573}{\kilo\joule\per\meter\cubed}$;
(b) $\approx\SI{40}{\joule\per\meter\cubed}$}
\resolucao{(a)
\[
u_B=\frac{B^{2}}{2\mu_0}=\frac{(1{,}2)^{2}}{2(4\pi\times10^{-7})}
=\frac{1{,}44}{\pot{2.513}{-6}}
\approx\pot{5.73}{5}\,\si{\joule\per\meter\cubed}.
\]
(b)
\[
u_E=\frac12\varepsilon_0E^{2}
=\frac12(\pot{8.85}{-12})(3\times10^{6})^{2}
\approx\SI{40}{\joule\per\meter\cubed}.
\]
\textbf{Razão: cerca de $14\,000$ vezes.} O campo elétrico usado é o limite de
ruptura do ar, e o magnético é típico de aço-silício --- ou seja, comparam-se
dois valores \emph{praticáveis}.\\
\textbf{Consequência tecnológica:} armazenar energia magneticamente é muito
mais denso do que eletrostaticamente. É por isso que transformadores e motores
operam com campos magnéticos, e por que capacitores só competem em energia
quando se usa dielétrico de altíssima constante e espessura nanométrica ---
como nos supercapacitores.}
\end{questao}

\begin{questao}[3]{}
Duas bobinas têm $L_1=\SI{50}{\milli\henry}$ e $L_2=\SI{200}{\milli\henry}$.
\begin{itensq}
\item Determine o valor máximo possível de $M$.
\item Se $M=\SI{80}{\milli\henry}$, determine $k$.
\item Interprete fisicamente $k=1$ e $k=0$.
\end{itensq}
\resposta{(a) $M_{\max}=\SI{100}{\milli\henry}$; (b) $k=0{,}8$}
\resolucao{(a) Como $k\le1$,
\[
M_{\max}=\sqrt{L_1L_2}=\sqrt{50\cdot200}=\sqrt{10\,000}
=\SI{100}{\milli\henry}.
\]
(b) $k=80/100=0{,}8$ --- acoplamento forte, típico de bobinas sobre um mesmo
núcleo ferromagnético com pequeno entreferro.\\
(c) $k=1$ significa \textbf{acoplamento perfeito}: \emph{todo} o fluxo produzido
por uma bobina concatena a outra, sem dispersão. É o transformador ideal ---
inatingível, mas aproximado em núcleos toroidais fechados de alta
permeabilidade.\\
$k=0$ significa \textbf{desacoplamento total}: nenhuma influência mútua. Obtém-se
afastando as bobinas ou dispondo seus eixos perpendicularmente, de modo que o
fluxo de uma atravesse a outra tangencialmente. É a razão de indutores vizinhos
em placas serem montados em orientações ortogonais.}
\end{questao}

\begin{questao}[3]{}
Um indutor com núcleo apresenta $L=\SI{100}{\micro\henry}$ até
$\SI{3}{\ampere}$; acima disso, a saturação reduz $L$ a
$\SI{20}{\micro\henry}$. Ele opera com $\SI{12}{\volt}$ aplicados durante
$\SI{10}{\micro\second}$ a partir de $\SI{2}{\ampere}$.
\begin{itensq}
\item Determine a corrente ao final do pulso, ignorando a saturação.
\item Refaça considerando a saturação.
\item Comente a consequência prática.
\end{itensq}
\resposta{(a) $\SI{3.2}{\ampere}$; (b) $\SI{5}{\ampere}$; (c) risco de
destruição do interruptor}
\resolucao{(a) Com $L$ constante,
\[
\Delta i=\frac{v\,\Delta t}{L}=\frac{12\cdot10^{-5}}{10^{-4}}
=\SI{1.2}{\ampere}
\;\Longrightarrow\;
i_{final}=2+1{,}2=\SI{3.2}{\ampere}.
\]
(b) A corrente sobe de $2$ a $\SI{3}{\ampere}$ com
$L=\SI{100}{\micro\henry}$, o que consome
\[
\Delta t_1=\frac{L\,\Delta i}{v}=\frac{10^{-4}\cdot1}{12}
=\SI{8.33}{\micro\second}.
\]
Restam $\SI{1.67}{\micro\second}$, agora com $L=\SI{20}{\micro\henry}$:
\[
\Delta i_2=\frac{12\cdot\pot{1.67}{-6}}{\pot{2}{-5}}=\SI{1}{\ampere}
\;\Longrightarrow\;
i_{final}=\SI{5}{\ampere}.
\]
(c) A corrente real é $56\%$ maior que a prevista pelo modelo linear, e a
subida final é \emph{cinco vezes} mais íngreme. Em um conversor chaveado, isso
significa que o transistor pode ver o dobro da corrente esperada em uma fração
do tempo --- destruição por sobrecorrente antes que qualquer proteção lenta
atue.\\
\textbf{Regra de projeto:} dimensione o indutor para que a corrente de pico
fique confortavelmente abaixo do joelho de saturação, e verifique
$L$ \emph{na corrente de operação}, nunca no valor de catálogo medido a
corrente nula.}
\end{questao}

\begin{questao}[3]{}
Uma fonte de \textbf{corrente} de $\SI{2}{\ampere}$ alimenta um indutor de
$\SI{0.5}{\henry}$ em paralelo com $R=\SI{10}{\ohm}$, em regime permanente.
\begin{itensq}
\item Determine a corrente no indutor e no resistor.
\item A fonte é subitamente desligada (aberta). Determine a tensão inicial
      sobre o resistor.
\item Determine o tempo de extinção.
\end{itensq}
\resposta{(a) $\SI{2}{\ampere}$ no indutor, $\SI{0}{\ampere}$ no resistor;
(b) $\SI{-20}{\volt}$; (c) $\tau=\SI{50}{\milli\second}$}
\resolucao{(a) Em regime CC o indutor ideal é um \textbf{curto}: ele desvia
toda a corrente, e o resistor em paralelo fica com tensão nula --- logo
corrente nula. Todos os $\SI{2}{\ampere}$ passam por $L$.\\
(b) Ao abrir a fonte, a corrente do indutor \emph{não pode saltar}: continua
$\SI{2}{\ampere}$, e o único caminho disponível é o resistor. A corrente agora
o percorre em sentido oposto ao anterior:
\[
v_R=-RI_0=-10(2)=\SI{-20}{\volt}.
\]
(c) $\tau=L/R=0{,}5/10=\SI{50}{\milli\second}$; a extinção é praticamente
completa em $5\tau=\SI{250}{\milli\second}$.\\
\textbf{Observação:} o resistor, que parecia inútil em regime (corrente nula),
é justamente o que evita a sobretensão destrutiva na abertura. Ele é um
\emph{caminho de escape} projetado --- exatamente a função do diodo de roda
livre, aqui desempenhada por um resistor.}
\end{questao}

\blocodesafio

\begin{questao}[4]{}
\textbf{(Dedução --- resposta do circuito RL.)} Um circuito RL série com fonte
$V$ tem a chave fechada em $t=0$, com $i(0)=0$.
\begin{itensq}
\item Escreva a equação diferencial e resolva-a por separação de variáveis.
\item Identifique $\tau$ e o valor final.
\item Obtenha $v_L(t)$ e verifique a LKT em $t=0^{+}$ e $t\to\infty$.
\end{itensq}
\resposta{$i(t)=\dfrac{V}{R}\left(1-e^{-Rt/L}\right)$;
$\tau=L/R$; $v_L=Ve^{-Rt/L}$}
\resolucao{(a) Pela LKT,
\[
V=Ri+L\frac{\mathrm{d}i}{\mathrm{d}t}
\;\Longrightarrow\;
\frac{\mathrm{d}i}{V/R-i}=\frac{R}{L}\,\mathrm{d}t.
\]
Integrando de $0$ a $t$ (com $i(0)=0$):
\[
-\ln\!\left(\frac{V/R-i}{V/R}\right)=\frac{R}{L}t
\;\Longrightarrow\;
i(t)=\frac{V}{R}\left(1-e^{-Rt/L}\right).
\]
(b) O expoente deve ser adimensional, o que identifica
$\tau=L/R$; e $i(\infty)=V/R$, coerente com o indutor tornando-se um curto.\\
(c)
\[
v_L=L\frac{\mathrm{d}i}{\mathrm{d}t}
=L\cdot\frac{V}{R}\cdot\frac{R}{L}e^{-Rt/L}=V\,e^{-Rt/L}.
\]
\emph{Em $t=0^{+}$:} $v_L=V$ e $Ri=0$; soma $=V$ \checkmark.\\
\emph{Em $t\to\infty$:} $v_L=0$ e $Ri=V$; soma $=V$ \checkmark.\\
\textbf{Estrutura geral:} a solução é (regime permanente) $+$ (transitório
exponencial). Essa decomposição vale para \emph{qualquer} circuito de primeira
ordem e dispensa refazer a integração a cada problema.}
\end{questao}

\begin{questao}[4]{}
\textbf{(Integração --- energia na descarga.)} Um indutor com corrente inicial
$I_0$ é descarregado através de $R$.
\begin{itensq}
\item Escreva $i(t)$ e a potência instantânea no resistor.
\item Integre para obter a energia total dissipada.
\item Mostre que o resultado independe de $R$ e interprete.
\end{itensq}
\resposta{$W=\frac12LI_0^{2}$, independente de $R$}
\resolucao{(a) Sem fonte, $Ri+L\,\mathrm{d}i/\mathrm{d}t=0$, cuja solução é
$i(t)=I_0e^{-t/\tau}$ com $\tau=L/R$. Então
\[
p_R(t)=Ri^{2}=RI_0^{2}e^{-2t/\tau}.
\]
(b)
\[
W=\int_0^{\infty}RI_0^{2}e^{-2t/\tau}\,\mathrm{d}t
=RI_0^{2}\left[-\frac{\tau}{2}e^{-2t/\tau}\right]_0^{\infty}
=\frac{RI_0^{2}\tau}{2}.
\]
Substituindo $\tau=L/R$:
\[
W=\frac{RI_0^{2}}{2}\cdot\frac{L}{R}=\frac12LI_0^{2}.
\]
(c) O $R$ \textbf{cancela} exatamente: ele aparece na potência
($\propto R$) e desaparece na duração ($\tau\propto1/R$). Aumentar $R$ dez
vezes decuplica a potência inicial e divide o tempo por dez --- a área sob a
curva não muda.\\
\textbf{Significado:} a energia dissipada é determinada pelo \emph{estado
inicial} do indutor, não pelo caminho de descarga. É a mesma lógica que faz um
capacitor descarregado entregar sempre $\frac12CV_0^{2}$, e confirma por
integração o argumento de conservação usado no bloco anterior.}
\end{questao}

\begin{questao}[4]{}
\textbf{(Projeto de proteção.)} Uma bobina de relé com $L=\SI{0.2}{\henry}$ e
$r=\SI{100}{\ohm}$ opera em $\SI{24}{\volt}$. O transistor de comando suporta
no máximo $\SI{60}{\volt}$.
\begin{itensq}
\item Determine a corrente de operação e a energia armazenada.
\item Dimensione um grampeador resistivo (diodo em série com $R_g$) que
      respeite o limite do transistor.
\item Compare o tempo de extinção com o do diodo simples.
\end{itensq}
\resposta{(a) $\SI{0.24}{\ampere}$, $\SI{5.76}{\milli\joule}$;
(b) $R_g\le\SI{150}{\ohm}$; (c) $\SI{0.8}{\milli\second}$ contra
$\SI{2}{\milli\second}$}
\resolucao{(a) $I=24/100=\SI{0.24}{\ampere}$;
$W=\frac12(0{,}2)(0{,}24)^{2}=\SI{5.76}{\milli\joule}$.\\
(b) Na abertura, a corrente $\SI{0.24}{\ampere}$ é forçada pelo diodo e por
$R_g$. A tensão no coletor sobe para $V+R_gI$ (desprezando o diodo). Exigindo
$\le\SI{60}{\volt}$:
\[
24+R_g(0{,}24)\le60
\;\Longrightarrow\;
R_g\le\frac{36}{0{,}24}=\SI{150}{\ohm}.
\]
Adote $R_g=\SI{150}{\ohm}$ e verifique a dissipação: o pulso entrega
$\SI{5.76}{\milli\joule}$ e, a $\SI{10}{\hertz}$ de comutação, isso são
$\SI{57.6}{\milli\watt}$ médios --- um resistor de
$\frac14\,\si{\watt}$ basta, desde que suporte o pico de
$0{,}24^{2}(150)=\SI{8.6}{\watt}$ instantâneos.\\
(c) \emph{Com grampeador:} $\tau=L/(r+R_g)=0{,}2/250=\SI{0.8}{\milli\second}$.\\
\emph{Com diodo simples:} $\tau=L/r=0{,}2/100=\SI{2}{\milli\second}$.\\
\textbf{Compromisso quantificado:} o grampeador reduz o tempo de desarme em
$60\%$, ao custo de submeter o transistor a $\SI{60}{\volt}$ em vez de
$\SI{24.7}{\volt}$. Quanto maior $R_g$, mais rápido o relé desarma e maior a
tensão no interruptor --- o projeto consiste em usar toda a margem do
transistor, e nada além dela.}
\end{questao}

\begin{questao}[4]{}
\textbf{(Projeto de armazenamento.)} Deseja-se um indutor que armazene
$\SI{5}{\joule}$ conduzindo $\SI{10}{\ampere}$.
\begin{itensq}
\item Determine $L$.
\item Para um núcleo toroidal com $A=\SI{10}{\centi\meter\squared}$,
      $\ell=\SI{0.25}{\meter}$ e $\mu_r=2000$, determine o número de espiras.
\item Verifique a densidade de fluxo e avalie a viabilidade.
\end{itensq}
\resposta{(a) $L=\SI{0.1}{\henry}$; (b) $N\approx35$ espiras;
(c) $B\approx\SI{3.5}{\tesla}$ --- \textbf{inviável}, o núcleo satura}
\resolucao{(a) $W=\frac12LI^{2}\Rightarrow
L=\dfrac{2W}{I^{2}}=\dfrac{10}{100}=\SI{0.1}{\henry}$.\\
(b) De $L=\dfrac{\mu_r\mu_0N^{2}A}{\ell}$:
\[
N=\sqrt{\frac{L\ell}{\mu_r\mu_0A}}
=\sqrt{\frac{0{,}1\cdot0{,}25}{2000(4\pi\times10^{-7})(10^{-3})}}
=\sqrt{9{,}95}\approx35\ \text{espiras}.
\]
(c) O fluxo é
\[
B=\frac{\mu_r\mu_0NI}{\ell}
=\frac{2000(4\pi\times10^{-7})(35)(10)}{0{,}25}\approx\SI{3.5}{\tesla}.
\]
Materiais ferromagnéticos comuns saturam entre $1{,}2$ e
$\SI{1.8}{\tesla}$. O projeto é \textbf{inviável}: muito antes de atingir
$\SI{10}{\ampere}$, o núcleo satura, $\mu_r$ despenca e $L$ some.\\
\textbf{Diagnóstico e correção.} A energia armazenada em material
ferromagnético é limitada justamente porque $\mu_r$ alto significa $B$ alto para
pouca corrente. A solução clássica é introduzir um \textbf{entreferro}: ele
reduz a permeabilidade efetiva (exigindo mais espiras), mas permite corrente
muito maior antes da saturação --- e a maior parte da energia passa a residir
no entreferro, onde $u=B^{2}/2\mu_0$ é enorme.\\
\textbf{Lição:} calcular $L$ e $N$ não basta. Sem verificar $B$, obtém-se um
projeto aritmeticamente correto e fisicamente impossível.}
\end{questao}

\begin{questao}[4]{}
\textbf{(Projeto de constante de tempo.)} Deseja-se $\tau=\SI{20}{\milli\second}$
com $R=\SI{50}{\ohm}$.
\begin{itensq}
\item Determine $L$.
\item Avalie a viabilidade física dessa indutância.
\item Proponha duas alternativas para obter o mesmo $\tau$ de forma prática.
\end{itensq}
\resposta{(a) $L=\SI{1}{\henry}$; (c) aumentar $R$, ou usar RC equivalente}
\resolucao{(a) $L=\tau R=0{,}020(50)=\SI{1}{\henry}$.\\
(b) \textbf{Um henry é um componente grande.} Exige núcleo ferromagnético e
centenas de espiras; a resistência do próprio enrolamento facilmente atinge
dezenas de ohms --- comparável aos $\SI{50}{\ohm}$ do projeto, o que
\emph{altera} o $\tau$ pretendido. Some-se a isso massa, custo, sensibilidade a
campos externos e saturação.\\
(c) \textbf{Alternativa 1 --- elevar $R$.} Com $R=\SI{500}{\ohm}$,
$L=\SI{10}{\henry}$... pior. O caminho correto é o inverso: fixar um $L$
razoável e ajustar $R$. Com $L=\SI{10}{\milli\henry}$, obter
$\tau=\SI{20}{\milli\second}$ exigiria $R=\SI{0.5}{\ohm}$ --- baixo demais para
ser prático em circuitos de sinal.\\
\textbf{Alternativa 2 --- usar RC.} Um capacitor de
$\SI{20}{\micro\farad}$ com $\SI{1}{\kilo\ohm}$ dá $\tau=\SI{20}{\milli\second}$
em um componente barato, pequeno e estável.\\
\textbf{Conclusão de engenharia:} para constantes de tempo da ordem de
milissegundos ou mais, \textbf{RC vence RL} quase sempre. Indutores são
preferíveis quando se precisa de armazenamento de energia com corrente elevada
ou de comportamento em frequência --- não para gerar atrasos.}
\end{questao}

\begin{questao}[4]{}
\textbf{(Demonstração --- acoplamento em série.)} Duas bobinas acopladas em
série conduzem a mesma corrente $i$.
\begin{itensq}
\item Deduza $L_{eq}=L_1+L_2\pm2M$ a partir das tensões induzidas.
\item Demonstre que $M\le\sqrt{L_1L_2}$, usando o fato de a energia ser não
      negativa.
\item Explique o significado da convenção dos pontos.
\end{itensq}
\resposta{(b) da condição $W\ge0$ para toda corrente segue
$M^{2}\le L_1L_2$}
\resolucao{(a) Cada bobina sofre autoindução \emph{e} indução mútua:
\[
v_1=L_1\frac{\mathrm{d}i}{\mathrm{d}t}\pm M\frac{\mathrm{d}i}{\mathrm{d}t},
\qquad
v_2=L_2\frac{\mathrm{d}i}{\mathrm{d}t}\pm M\frac{\mathrm{d}i}{\mathrm{d}t}.
\]
Somando e comparando com
$v=L_{eq}\,\mathrm{d}i/\mathrm{d}t$:
\[
L_{eq}=L_1+L_2\pm2M.
\]
(b) Com correntes independentes $i_1$ e $i_2$, a energia do par é
\[
W=\frac12L_1i_1^{2}+\frac12L_2i_2^{2}+Mi_1i_2.
\]
Como um sistema passivo não pode armazenar energia negativa, $W\ge0$ para
\emph{quaisquer} $i_1,i_2$. Tomando $i_1=x$ e $i_2=-1$:
\[
\frac12L_1x^{2}-Mx+\frac12L_2\ge0\quad\forall x,
\]
o que exige discriminante não positivo:
\[
M^{2}-L_1L_2\le0
\;\Longrightarrow\;
M\le\sqrt{L_1L_2}
\;\Longrightarrow\;
k=\frac{M}{\sqrt{L_1L_2}}\le1.
\]
(c) Os pontos marcam terminais \textbf{magneticamente correspondentes}:
correntes que entram pelos terminais pontuados produzem fluxos que se
\emph{reforçam}. Se ambas as correntes entram (ou ambas saem) pelos pontos, o
termo mútuo é positivo; caso contrário, negativo.\\
\textbf{Por que a convenção é necessária:} o esquema elétrico não informa o
sentido de enrolamento das bobinas, que é uma propriedade \emph{física} da
construção. Sem os pontos, $L_{eq}$ seria ambíguo entre
$\SI{19}{\milli\henry}$ e $\SI{7}{\milli\henry}$ --- e, em um transformador, a
polaridade da saída ficaria indeterminada.}
\end{questao}

\begin{questao}[4]{}
\textbf{(Indutor real em regime.)} Um indutor com $L=\SI{0.5}{\henry}$ e
$r=\SI{2}{\ohm}$ é alimentado por $\SI{10}{\volt}$ CC.
\begin{itensq}
\item Determine a corrente, a energia armazenada e a potência dissipada em
      regime.
\item Determine o tempo necessário para atingir $99\%$ da corrente final.
\item Discuta o "custo de manutenção" da energia armazenada.
\end{itensq}
\resposta{(a) $\SI{5}{\ampere}$, $\SI{6.25}{\joule}$, $\SI{50}{\watt}$;
(b) $\approx\SI{1.15}{\second}$}
\resolucao{(a) $I=10/2=\SI{5}{\ampere}$;
$W=\frac12(0{,}5)(25)=\SI{6.25}{\joule}$;
$P=rI^{2}=2(25)=\SI{50}{\watt}$.\\
(b) $\tau=L/r=0{,}5/2=\SI{0.25}{\second}$, e $99\%$ correspondem a
$t=\tau\ln100=0{,}25(4{,}605)=\SI{1.15}{\second}$.\\
(c) \textbf{O contraste é gritante.} Carregar o campo custou
$\SI{6.25}{\joule}$; mantê-lo custa $\SI{50}{\watt}$ \emph{continuamente}.
Em apenas $\SI{0.125}{\second}$ de operação, a dissipação já iguala toda a
energia armazenada; em um minuto, dissipam-se $\SI{3000}{\joule}$ --- $480$
vezes o conteúdo do campo.\\
\textbf{Figura de mérito:} a razão $W/P=L/(2r)=\tau/2$ tem dimensão de tempo e
mede a "qualidade" do indutor em CC. Aqui vale $\SI{0.125}{\second}$.\\
\textbf{Consequência:} indutores são péssimos \emph{reservatórios} de energia
para armazenamento prolongado --- ao contrário de baterias ou capacitores, que
não exigem corrente para manter o estado. Eles são excelentes para
\emph{transferir} energia em ciclos curtos, que é exatamente o que fazem em
conversores chaveados, onde o ciclo dura microssegundos.}
\end{questao}

\begin{questao}[4]{}
\textbf{(Conversor chaveado.)} Um conversor opera a $\SI{50}{\kilo\hertz}$ com
razão cíclica $D=0{,}5$, aplicando $\SI{12}{\volt}$ sobre um indutor de
$\SI{100}{\micro\henry}$ durante a condução.
\begin{itensq}
\item Determine a ondulação de corrente pico a pico.
\item Determine a corrente de pico se a componente contínua for
      $\SI{3}{\ampere}$.
\item O indutor satura em $\SI{3.5}{\ampere}$. Proponha e quantifique duas
      correções.
\end{itensq}
\resposta{(a) $\Delta i=\SI{1.2}{\ampere}$; (b) $I_{pico}=\SI{3.6}{\ampere}$;
(c) satura --- elevar $L$ ou a frequência}
\resolucao{(a) Durante o tempo de condução
$t_{on}=D/f=0{,}5/\pot{5}{4}=\SI{10}{\micro\second}$:
\[
\Delta i=\frac{V\,t_{on}}{L}=\frac{12(10^{-5})}{10^{-4}}=\SI{1.2}{\ampere}.
\]
(b) A ondulação se distribui simetricamente em torno da média:
\[
I_{pico}=3+\frac{1{,}2}{2}=\SI{3.6}{\ampere}\ >\ \SI{3.5}{\ampere}.
\]
\textbf{O indutor satura} --- e note que a corrente \emph{média}
($\SI{3}{\ampere}$) estava confortavelmente abaixo do limite. Dimensionar pela
média é o erro clássico.\\
(c) Como $\Delta i=\dfrac{VD}{fL}$, há dois caminhos:\\
\textbf{(i) Aumentar $L$.} Para $I_{pico}\le3{,}5$ é preciso
$\Delta i\le\SI{1}{\ampere}$, ou seja
\[
L\ge\frac{12(10^{-5})}{1}=\SI{120}{\micro\henry}.
\]
Custo: indutor maior e mais caro.\\
\textbf{(ii) Aumentar a frequência.} Mantendo
$L=\SI{100}{\micro\henry}$, exige-se
\[
f\ge\frac{12(0{,}5)}{10^{-4}(1)}=\SI{60}{\kilo\hertz}.
\]
Custo: mais perdas de comutação no transistor e maior interferência
eletromagnética.\\
\textbf{Compromisso:} $\Delta i\propto1/(fL)$ --- o produto $fL$ é o que
importa. Elevar a frequência permite indutores menores, e é exatamente essa a
tendência que tornou as fontes chaveadas modernas tão compactas.}
\end{questao}

\begin{questao}[4]{}
\textbf{(Demonstração --- continuidade da corrente.)} Prove que a corrente em
um indutor ideal não pode sofrer descontinuidade, usando um argumento
energético.
\begin{itensq}
\item Suponha um salto de $I_1$ para $I_2$ em intervalo $\Delta t\to0$ e
      analise a tensão.
\item Analise a potência e a energia envolvidas.
\item Explique por que capacitores em paralelo apresentam o problema dual.
\end{itensq}
\resposta{Um salto exigiria tensão e potência infinitas; a energia é finita,
mas a taxa de transferência não pode ser}
\resolucao{(a) Um salto implica
\[
v=L\lim_{\Delta t\to0}\frac{I_2-I_1}{\Delta t}\to\infty.
\]
Nenhuma fonte real fornece tensão infinita.\\
(b) A potência $p=vi$ também diverge. A \emph{energia} necessária é finita
---
\[
\Delta W=\frac12L\left(I_2^{2}-I_1^{2}\right),
\]
--- mas ela teria de ser transferida em tempo nulo, exigindo potência
infinita. \textbf{É a taxa, não a quantidade, que torna o salto impossível.}\\
(c) \textbf{Problema dual:} dois capacitores com tensões diferentes ligados em
paralelo exigiriam salto de \emph{tensão}, o que demandaria corrente infinita
($i=C\,\mathrm{d}v/\mathrm{d}t$). Na prática, a resistência inevitável dos
condutores limita a corrente e dissipa a diferença de energia --- e demonstra-se
que a perda independe do valor dessa resistência, resultado análogo ao da
descarga do indutor.\\
\textbf{Regra prática:} nunca ligue em série indutores com correntes distintas,
nem em paralelo capacitores com tensões distintas. Em ambos os casos, o modelo
ideal falha e o circuito real encontra uma saída --- arco elétrico no primeiro
caso, pico de corrente no segundo.}
\end{questao}

\begin{questao}[4]{}
\textbf{(Modelagem --- sensor indutivo.)} Um sensor de proximidade usa uma
bobina cuja indutância varia com a distância $x$ a um alvo ferromagnético,
segundo $L(x)=L_0\left(1+\dfrac{a}{x}\right)$, com
$L_0=\SI{10}{\milli\henry}$ e $a=\SI{1}{\milli\meter}$.
\begin{itensq}
\item Determine $L$ para $x=\SI{1}{\milli\meter}$ e $x=\SI{5}{\milli\meter}$.
\item Obtenha a sensibilidade $\mathrm{d}L/\mathrm{d}x$ e avalie nos dois
      pontos.
\item Discuta a consequência para o projeto do sensor.
\end{itensq}
\resposta{(a) $\SI{20}{\milli\henry}$ e $\SI{12}{\milli\henry}$;
(b) $-\SI{10}{\milli\henry\per\milli\meter}$ e
$-\SI{0.4}{\milli\henry\per\milli\meter}$}
\resolucao{(a)
\[
L(1)=10(1+1)=\SI{20}{\milli\henry};
\qquad
L(5)=10(1+0{,}2)=\SI{12}{\milli\henry}.
\]
(b)
\[
\frac{\mathrm{d}L}{\mathrm{d}x}=-\frac{L_0a}{x^{2}}.
\]
Em $x=\SI{1}{\milli\meter}$:
$-10(1)/1=-\SI{10}{\milli\henry\per\milli\meter}$.\\
Em $x=\SI{5}{\milli\meter}$:
$-10(1)/25=-\SI{0.4}{\milli\henry\per\milli\meter}$.\\
A sensibilidade cai com $1/x^{2}$: é \textbf{25 vezes menor} a
$\SI{5}{\milli\meter}$.\\
(c) \textbf{Três consequências.} \emph{(i)} O sensor é fortemente
\textbf{não linear} --- a leitura bruta de $L$ não é proporcional a $x$, e a
conversão exige linearização ou tabela. \emph{(ii)} A \textbf{faixa útil} é
curta: além de poucos milímetros, a variação se confunde com a deriva térmica
da própria bobina. \emph{(iii)} A \textbf{resolução} depende da posição: perto
do alvo, décimos de micrômetro são detectáveis; longe, nem décimos de
milímetro.\\
\textbf{Solução usual:} operar sempre na região próxima e usar o sensor como
detector de \emph{limiar} (presença/ausência), em vez de medidor de distância
--- que é exatamente como os sensores indutivos industriais são
especificados.}
\end{questao}

% END BANCO COMPLEMENTAR

\gabarito[10.2 Indutores]
